% =====================================================================
% Chapter 5: Why the gap sits at 1031--1195 Hz and not at 1200--1342 Hz
% Theory session, dispersion branch (Nov 2026)
% Requires: amsmath, amssymb
% Include in the main file with  % =====================================================================
% Chapter 5: Why the gap sits at 1031--1195 Hz and not at 1200--1342 Hz
% Theory session, dispersion branch (Nov 2026)
% Requires: amsmath, amssymb
% Include in the main file with  % =====================================================================
% Chapter 5: Why the gap sits at 1031--1195 Hz and not at 1200--1342 Hz
% Theory session, dispersion branch (Nov 2026)
% Requires: amsmath, amssymb
% Include in the main file with  % =====================================================================
% Chapter 5: Why the gap sits at 1031--1195 Hz and not at 1200--1342 Hz
% Theory session, dispersion branch (Nov 2026)
% Requires: amsmath, amssymb
% Include in the main file with  \include{chapter_05_gap_position}
% =====================================================================

\chapter{Why the gap sits at 1031--1195\,Hz and not at 1200--1342\,Hz}
\label{ch:gap_position}

The effective-medium estimate of Chapter~\ref{ch:local_resonance} predicts a
gap of $1200$--$1342\,\mathrm{Hz}$. The FE diagram shows a complete gap of
$1031$--$1195\,\mathrm{Hz}$. This chapter explains the difference. Every
equation is derived here, starting from Newton's law for the resonator, so the
chapter can be read on its own.

Unless stated otherwise, numbers come from the Reissner--Mindlin (RM)
plane-wave model, which agrees with FE to within $0.33\,\%$. The exact bare
frequency at $X$ is $1194.4\,\mathrm{Hz}$; earlier chapters round it to
$1190\,\mathrm{Hz}$. No conclusion there changes.

\begin{center}
\begin{tabular}{|l|p{0.62\textwidth}|}
\hline
Symbol & Meaning \\
\hline
$\omega$, $f$ & angular frequency and frequency of a wave (one dot of the diagram) \\
$\lambda = \omega^2$ & eigenvalue belonging to that dot \\
$\mathbf{k}$ & Bloch wavevector, the horizontal axis of the diagram \\
$\mathbf{x}$, $\mathbf{x}_0$ & position on the plate; position of the resonator (origin placed there, $\mathbf{x}_0 = 0$) \\
$a$, $A = a^2$ & cell size ($0.1\,\mathrm{m}$) and cell area \\
$\rho h$ & plate mass per area \\
$m$, $k_s$ & resonator mass and spring stiffness \\
$\lambda_r = \omega_r^2 = k_s/m$ & the resonator's own eigenvalue ($f_r = 1200\,\mathrm{Hz}$ unless stated) \\
$\mu = m/(\rho h A)$ & resonator mass over plate mass per cell, $0.25$ \\
$\beta = \lambda/\lambda_r = (f/f_r)^2$ & squared frequency relative to the resonator \\
$\varphi_n(\mathbf{x})$ & bare plate wave $n$, defined below \\
$\lambda_n^0$ & bare eigenvalue of wave $n$ (plate without resonator) \\
$\lambda_p$ & $\lambda_n^0$ of the one wave, or group of waves, kept in a simplified model \\
$q_n$ & amplitude of wave $n$ in the plate motion \\
$w_0$ & plate deflection at the resonator \\
$u$ & displacement of the resonator mass \\
$F$ & force the resonator exerts on the plate \\
$Z$, $C$ & force per deflection of the resonator; deflection per force of the plate \\
$S$ & sum over all waves except a chosen group, defined where it is used \\
$d$ & number of bare waves sharing the same $\lambda_n^0$ at a given $\mathbf{k}$ \\
$f_X$ & bare plate frequency at $X$, $1194.4\,\mathrm{Hz}$ \\
\hline
\end{tabular}
\end{center}

% ---------------------------------------------------------------------
\section{Recap: what the diagram contains}

A bare plate wave in the cell, for a given $\mathbf{k}$, is
\begin{equation}
  \varphi_n(\mathbf{x}) = \frac{e^{i(\mathbf{k}+\mathbf{G}_n)\cdot\mathbf{x}}}{\sqrt{\rho h A}} .
\end{equation}
The $\mathbf{G}_n = (2\pi/a)\,(\text{integer}, \text{integer})$ are the reciprocal
lattice vectors. Several waves fit at the same $\mathbf{k}$ because adding
$\mathbf{G}$ changes nothing from cell to cell. Each wave has its own bare
eigenvalue $\lambda_n^0$, set by its wavenumber $|\mathbf{k}+\mathbf{G}_n|$.

The factor $1/\sqrt{\rho h A}$ is chosen so that
$\int \rho h\,|\varphi_n|^2\,\mathrm{d}A = 1$ over the cell. Since
$|e^{i\dots}| = 1$,
\begin{equation}
  |\varphi_n(\mathbf{x}_0)|^2 = \frac{1}{\rho h A} \quad\text{for every wave } n.
  \label{eq:phi_squared}
\end{equation}

\begin{center}
\begin{tabular}{|l|l|l|l|}
\hline
Point & $\mathbf{k}$ & Bare waves there ($d$ of them) & Bare frequency \\
\hline
$\Gamma$ & $(0,0)$ & 1 wave, $\mathbf{G} = 0$ & $0\,\mathrm{Hz}$ \\
$\Gamma$ & $(0,0)$ & 4 waves, wavelength $a$ & $4715\,\mathrm{Hz}$ \\
$X$ & $(\pi/a, 0)$ & 2 waves, $(\pm\pi/a, 0)$ & $1194.4\,\mathrm{Hz}$ \\
$M$ & $(\pi/a, \pi/a)$ & 4 waves, $(\pm\pi/a, \pm\pi/a)$ & $2378.2\,\mathrm{Hz}$ \\
\hline
\end{tabular}
\end{center}

A \emph{complete gap} is a frequency interval with no dot at any $\mathbf{k}$.
It runs from the highest point of band 1 to the lowest point of band 2.

% ---------------------------------------------------------------------
\section{The resonator as a mass that depends on frequency}

Newton's law for the resonator mass: the spring force on it is
$k_s(w_0 - u)$, and for harmonic motion $\ddot{u} = -\omega^2 u$, so
\begin{align}
  -m\omega^2 u &= k_s(w_0 - u) \\
  u\,(k_s - m\omega^2) &= k_s w_0 \\
  u &= w_0\,\frac{k_s}{k_s - m\omega^2} .
\end{align}
Divide numerator and denominator by $k_s$. Since
$m\omega^2/k_s = \omega^2/\omega_r^2 = \beta$,
\begin{equation}
  u = \frac{w_0}{1 - \beta} .
  \label{eq:u_over_w}
\end{equation}
The force on the plate is equal and opposite to the force on the mass (Newton's
third law). The force on the mass is $m\ddot{u} = -m\omega^2 u$, so
\begin{equation}
  F = m\omega^2 u = \frac{m\,\omega^2 w_0}{1 - \beta} .
\end{equation}
A \emph{rigid} mass $M_\mathrm{rigid}$ glued to the plate moves with $w_0$, and
by the same reasoning exerts $F = M_\mathrm{rigid}\,\omega^2 w_0$. The two forces
have the same form. To the plate, the resonator therefore looks like a rigid
mass of size
\begin{equation}
  \boxed{\;m_\mathrm{eff} = \frac{m}{1-\beta}\;}
  \qquad\text{and, divided by } \rho h A, \qquad
  \boxed{\;\mu_\mathrm{eff} = \frac{\mu}{1-\beta}\;}
\end{equation}
Since $1 - \beta = (\lambda_r - \lambda)/\lambda_r$, an equivalent form is
\begin{equation}
  \mu_\mathrm{eff} = \frac{\mu\,\lambda_r}{\lambda_r - \lambda} .
  \label{eq:mu_eff_lambda}
\end{equation}
\begin{itemize}
  \item $f$ far below $f_r$ ($\beta \to 0$): $m_\mathrm{eff} \to m$; the spring
        is effectively rigid.
  \item $f$ rising towards $f_r$ ($\beta \to 1$): $m_\mathrm{eff}$ becomes very
        large and positive.
  \item $f$ above $f_r$ ($\beta > 1$): $m_\mathrm{eff}$ is negative.
\end{itemize}
In short, with $\lambda = \omega^2$,
\begin{equation}
  F = m_\mathrm{eff}\,\lambda\, w_0 .
  \label{eq:force}
\end{equation}

% ---------------------------------------------------------------------
\section{The plate driven by that force}

The plate with a point force $F$ at $\mathbf{x}_0$ obeys, with $\mathcal{L}$ the
plate stiffness operator,
\begin{equation}
  \mathcal{L} w - \rho h\,\lambda\, w = F\,\delta(\mathbf{x} - \mathbf{x}_0) .
\end{equation}
Expand the plate motion in bare waves, $w(\mathbf{x}) = \sum_j q_j\,\varphi_j(\mathbf{x})$,
multiply by $\varphi_n^*$ and integrate over the cell:
\begin{itemize}
  \item the stiffness term gives $\lambda_n^0\, q_n$ ($\varphi_n$ is a bare wave
        and the waves are orthogonal);
  \item the inertia term gives $\lambda\, q_n$ (normalization);
  \item on the right, the delta function picks the value at $\mathbf{x}_0$,
        giving $F\,\varphi_n(\mathbf{x}_0)^*$.
\end{itemize}
Therefore, for every wave $n$,
\begin{equation}
  (\lambda_n^0 - \lambda)\, q_n = F\,\varphi_n(\mathbf{x}_0)^* .
  \label{eq:modal}
\end{equation}
This says how strongly each wave responds to the point force.

\emph{Remark for Reissner--Mindlin:} there $\varphi_n$ also has rotation
components, and Eq.~\eqref{eq:modal} holds with the deflection part of
$\varphi_n$. The structure is identical, and the RM numbers include this.

% ---------------------------------------------------------------------
\section{Closing the loop: the plane-wave expansion condition}

Eq.~\eqref{eq:modal} gives the amplitude of each wave,
\begin{equation}
  q_n = \frac{F\,\varphi_n(\mathbf{x}_0)^*}{\lambda_n^0 - \lambda} ,
  \label{eq:qn}
\end{equation}
and the deflection at the resonator is the sum of all waves there,
\begin{equation}
  w_0 = \sum_n \varphi_n(\mathbf{x}_0)\, q_n .
  \label{eq:w0_sum}
\end{equation}

\paragraph{Insert \eqref{eq:qn} into \eqref{eq:w0_sum}.}
\begin{equation}
  w_0 = \sum_n \varphi_n(\mathbf{x}_0)\cdot\frac{F\,\varphi_n(\mathbf{x}_0)^*}{\lambda_n^0 - \lambda} .
\end{equation}
$F$ does not depend on $n$, so it comes out of the sum, and
$\varphi_n\varphi_n^* = |\varphi_n|^2$:
\begin{equation}
  w_0 = F\,\sum_n \frac{|\varphi_n(\mathbf{x}_0)|^2}{\lambda_n^0 - \lambda} .
\end{equation}
Define the sum as the \emph{compliance} $C$, the deflection per unit force:
\begin{equation}
  C = \sum_n \frac{|\varphi_n(\mathbf{x}_0)|^2}{\lambda_n^0 - \lambda} ,
  \qquad
  w_0 = C\,F .
  \label{eq:compliance}
\end{equation}
Likewise call the factor in Eq.~\eqref{eq:force} $Z = m_\mathrm{eff}\,\lambda$,
the force per unit deflection:
\begin{equation}
  F = Z\, w_0 .
  \label{eq:Z}
\end{equation}

\paragraph{Insert \eqref{eq:Z} into \eqref{eq:compliance} and divide by $w_0$.}
\begin{equation}
  w_0 = C\, Z\, w_0
  \quad\Longrightarrow\quad
  1 = Z\, C \qquad (w_0 \neq 0).
  \label{eq:ZC}
\end{equation}
The deflection creates a force, and that force creates exactly the same
deflection: the condition for a wave that sustains itself.

\paragraph{Rewriting $1 = ZC$ with mass ratios.}
Insert the definitions of $Z$ and $C$:
\begin{equation}
  1 = m_\mathrm{eff}\,\lambda\,\sum_n \frac{|\varphi_n(\mathbf{x}_0)|^2}{\lambda_n^0 - \lambda} .
\end{equation}
By Eq.~\eqref{eq:phi_squared}, $|\varphi_n(\mathbf{x}_0)|^2 = 1/(\rho h A)$ for
every $n$, so it comes out of the sum:
\begin{equation}
  1 = m_\mathrm{eff}\,\lambda\,\frac{1}{\rho h A}\,\sum_n \frac{1}{\lambda_n^0 - \lambda} .
\end{equation}
Group $m_\mathrm{eff}/(\rho h A) = \mu_\mathrm{eff}$ and move $\lambda$ into the
sum:
\begin{equation}
  \boxed{\;1 = \mu_\mathrm{eff}\,\sum_n \frac{\lambda}{\lambda_n^0 - \lambda}\;}
  \label{eq:pwe}
\end{equation}
This is the \emph{plane-wave expansion (PWE) condition}. At a given $\mathbf{k}$,
every $\lambda$ that satisfies it is one dot of the diagram. ``Kirchhoff'' or
``RM'' only decides which values $\lambda_n^0$ go into the sum; a ``Kirchhoff
compliance'' is $C$ evaluated with Kirchhoff values. In the notes the same
condition is written $1 = Z/A\cdot\sum_G g$, with
$g = 1/(\rho h(\omega_n^2 - \omega^2))$.

\paragraph{What if $w_0 = 0$?}
Then $F = 0$ by Eq.~\eqref{eq:Z}, and Eq.~\eqref{eq:modal} becomes
$(\lambda_n^0 - \lambda)\, q_n = 0$. A nonzero wave requires
$\lambda = \lambda_n^0$. A mode with a node at the resonator keeps its bare
frequency. Such modes do not appear in Eq.~\eqref{eq:pwe}, because we divided
by $w_0$, and must be added by hand.

% ---------------------------------------------------------------------
\section{Why a wave acts like a spring or like a mass}

Take one wave and write $g = F\,\varphi_n(\mathbf{x}_0)^*$ for the force it
feels. Eq.~\eqref{eq:qn} reads
\begin{equation}
  q_n = \frac{g}{\lambda_n^0 - \lambda} .
\end{equation}
Compare with the two elementary elements (unit mass, deflection $s$):
\begin{itemize}
  \item a \emph{spring} of stiffness $\kappa_s$: $\kappa_s\, s = g$, so
        $s = g/\kappa_s$. The deflection is along the force and does not depend
        on frequency.
  \item a \emph{mass} of $1$: $-\lambda\, s = g$, so $s = -g/\lambda$. The
        deflection is against the force.
\end{itemize}
Now the denominator decides:
\begin{itemize}
  \item \textbf{Wave above the driving frequency,} $\lambda_n^0 \gg \lambda$:
        $\lambda_n^0 - \lambda \approx \lambda_n^0 > 0$, so
        $q_n \approx g/\lambda_n^0$. This is exactly the spring formula with
        $\kappa_s = \lambda_n^0$.
  \item \textbf{Wave below the driving frequency,} $\lambda_n^0 \ll \lambda$:
        $\lambda_n^0 - \lambda \approx -\lambda < 0$, so
        $q_n \approx -g/\lambda$. This is exactly the mass formula.
\end{itemize}
Physically, each wave is an oscillator with a stiffness part $\lambda_n^0$ and
an inertia part $-\lambda$; the larger one decides its behaviour. For
$\lambda_n^0 = 0$ (no stiffness at all) the formula is exactly a mass.

% ---------------------------------------------------------------------
\section{One wave: where 1200--1342\,Hz comes from}

Keep only one term in Eq.~\eqref{eq:pwe}, with bare eigenvalue $\lambda_p$, and
use Eq.~\eqref{eq:mu_eff_lambda}:
\begin{equation}
  1 = \frac{\mu\,\lambda_r}{\lambda_r - \lambda}\cdot\frac{\lambda}{\lambda_p - \lambda} .
\end{equation}
Multiply both sides by $(\lambda_r - \lambda)(\lambda_p - \lambda)$:
\begin{equation}
  \boxed{\;(\lambda_p - \lambda)(\lambda_r - \lambda) = \mu\,\lambda_r\,\lambda\;}
  \label{eq:one_wave}
\end{equation}

\paragraph{Same equation as the determinant of Chapter~\ref{ch:local_resonance}.}
The two-level model with matrix
$H = \bigl(\begin{smallmatrix}\lambda_p+\mu\lambda_r & -\sqrt{\mu}\lambda_r\\ -\sqrt{\mu}\lambda_r & \lambda_r\end{smallmatrix}\bigr)$
gave $(\lambda_p + \mu\lambda_r - \lambda)(\lambda_r - \lambda) - \mu\lambda_r^2 = 0$.
Group the first bracket as $(\lambda_p - \lambda) + \mu\lambda_r$:
\begin{align}
  \bigl[(\lambda_p - \lambda) + \mu\lambda_r\bigr](\lambda_r - \lambda) - \mu\lambda_r^2 &= 0 \\
  (\lambda_p - \lambda)(\lambda_r - \lambda) + \mu\lambda_r(\lambda_r - \lambda) - \mu\lambda_r^2 &= 0 \\
  (\lambda_p - \lambda)(\lambda_r - \lambda) + \mu\lambda_r^2 - \mu\lambda_r\lambda - \mu\lambda_r^2 &= 0 \\
  (\lambda_p - \lambda)(\lambda_r - \lambda) - \mu\lambda_r\lambda &= 0 ,
\end{align}
which is Eq.~\eqref{eq:one_wave}. This form is useful because the left side
contains only the two uncoupled systems (it vanishes at $\lambda = \lambda_p$
and $\lambda = \lambda_r$), and the right side contains the coupling. With
$\mu = 0$ the uncoupled solutions come back.

\paragraph{The two bands.}
For each $\lambda_p$, Eq.~\eqref{eq:one_wave} has two solutions for $\lambda$.
Multiplying out gives a quadratic,
\begin{equation}
  \lambda^2 - \bigl(\lambda_p + \lambda_r(1+\mu)\bigr)\lambda + \lambda_p\lambda_r = 0 .
  \label{eq:quadratic}
\end{equation}
Call its two roots $\lambda_1 < \lambda_2$; they are the two bands at this $\mathbf{k}$. Writing the quadratic as $(\lambda - \lambda_1)(\lambda - \lambda_2) = \lambda^2 - (\lambda_1 + \lambda_2)\lambda + \lambda_1\lambda_2 = 0$ and comparing with Eq.~\eqref{eq:quadratic} gives $\lambda_1 + \lambda_2 = \lambda_p + \lambda_r(1+\mu)$ and $\lambda_1\lambda_2 = \lambda_p\lambda_r$. These two expressions are not the roots themselves, only their sum and their product.

\paragraph{Why two limits?}
The estimate pretends the plate is a uniform medium without a lattice. Then
$\lambda_p$ simply grows with the wavenumber, from $0$ at $\mathbf{k}=0$ towards
infinity. The bands as functions of $\lambda_p$, from Eq.~\eqref{eq:quadratic}:
\begin{center}
\begin{tabular}{|l|l|l|}
\hline
Bare plate frequency & Lower band & Upper band \\
\hline
$0\,\mathrm{Hz}$ & $0\,\mathrm{Hz}$ & $1342\,\mathrm{Hz}$ \\
$1000\,\mathrm{Hz}$ & $824\,\mathrm{Hz}$ & $1456\,\mathrm{Hz}$ \\
$2378\,\mathrm{Hz}$ & $1153\,\mathrm{Hz}$ & $2475\,\mathrm{Hz}$ \\
$10000\,\mathrm{Hz}$ & $1198\,\mathrm{Hz}$ & $10018\,\mathrm{Hz}$ \\
\hline
\end{tabular}
\end{center}
Both bands rise as $\lambda_p$ rises. The lower band is therefore \emph{highest}
at $\lambda_p \to \infty$ and the upper band \emph{lowest} at $\lambda_p = 0$.
These two ends bound the gap.

\paragraph{Limit $\lambda_p = 0$.}
Insert into Eq.~\eqref{eq:one_wave}:
$(-\lambda)(\lambda_r - \lambda) = \mu\lambda_r\lambda$. One solution is
$\lambda = 0$. Otherwise divide by $\lambda$:
\begin{align}
  -(\lambda_r - \lambda) &= \mu\lambda_r \\
  \lambda &= \lambda_r(1+\mu) \\
  f &= f_r\sqrt{1+\mu} = 1200\cdot\sqrt{1.25} = 1342\,\mathrm{Hz}.
\end{align}

\paragraph{Limit $\lambda_p \to \infty$.}
For very large $\lambda_p$ the sum of the roots is dominated by $\lambda_p$, so
the upper root is $\approx \lambda_p$. Then
\begin{equation}
  \text{lower root} = \frac{\text{product}}{\text{upper root}}
  \approx \frac{\lambda_p\lambda_r}{\lambda_p} = \lambda_r ,
\end{equation}
that is, $1200\,\mathrm{Hz}$.

\paragraph{The two hidden assumptions.}
\begin{itemize}
  \item \textbf{Assumption A:} $\lambda_p$ can grow without limit. In the
        lattice the zone ends at $M$, where the bare value stops at
        $2378\,\mathrm{Hz}$.
  \item \textbf{Assumption B:} one wave per $\mathbf{k}$ is enough. A point
        force excites all waves.
\end{itemize}

% ---------------------------------------------------------------------
\section{Several waves with the same bare frequency}

If $d$ waves share the same $\lambda_p$ and the others are ignored, the $d$ terms
in Eq.~\eqref{eq:pwe} are identical:
\begin{equation}
  1 = d\,\mu_\mathrm{eff}\,\frac{\lambda}{\lambda_p - \lambda} ,
\end{equation}
and the same multiplication gives
\begin{equation}
  \boxed{\;(\lambda_p - \lambda)(\lambda_r - \lambda) = d\,\mu\,\lambda_r\,\lambda\;}
  \label{eq:d_waves}
\end{equation}
This is Eq.~\eqref{eq:one_wave} with $\mu$ replaced by $d\mu$: $d$ waves act like
$d$ times the mass ratio. Of those $d$ waves, one combination has $w_0 \neq 0$;
the other $d-1$ combinations have a node at the resonator and stay at the bare
frequency.

\paragraph{Plate share of the motion.}
Let the coupling combination have amplitude $q$, normalized like the bare waves.
With $d$ waves adding at $\mathbf{x}_0$, $|w_0|^2 = q^2\,d/(\rho h A)$. The
kinetic energy in the plate is proportional to $q^2$, in the resonator to
$m\,u^2$. Using Eq.~\eqref{eq:u_over_w},
\begin{equation}
  m\,u^2 = \frac{m\, q^2\, d/(\rho h A)}{(1-\beta)^2} = q^2\,\frac{d\,\mu}{(1-\beta)^2} .
\end{equation}
Therefore
\begin{equation}
  \text{plate share} = \frac{q^2}{q^2 + m\,u^2} = \frac{1}{1 + d\mu/(1-\beta)^2} .
  \label{eq:plate_share}
\end{equation}
It is the same result as the eigenvector ratio of
Chapter~\ref{ch:local_resonance}, written in terms of $\beta$.

% ---------------------------------------------------------------------
\section{Upper edge: 1342 $\to$ 1273.5 $\to$ 1194.4\,Hz}

\paragraph{Step 1: stay at $\Gamma$, keep all waves.}
The sum in Eq.~\eqref{eq:pwe} has two parts:
\begin{itemize}
  \item the $\mathbf{G}=0$ wave has $\lambda_n^0 = 0$, and its term is
        $\lambda/(0-\lambda) = -1$: the mass-like part, the whole cell moving
        rigidly;
  \item all other waves lie far above ($4715\,\mathrm{Hz}$ and more). Their sum
        is given the name
        \begin{equation}
          S = \sum_{n \neq 0} \frac{\lambda}{\lambda_n^0 - \lambda} > 0 ,
        \end{equation}
        which is spring-like.
\end{itemize}
Eq.~\eqref{eq:pwe} becomes $1 = \mu_\mathrm{eff}(-1 + S)$. Insert
$\mu_\mathrm{eff} = \mu/(1-\beta)$ and multiply by $(1-\beta)$:
\begin{equation}
  1 - \beta = -\mu + \mu S
  \quad\Longrightarrow\quad
  \boxed{\;\beta = 1 + \mu - \mu S\;}
\end{equation}
$S$ itself depends on $\beta$, so the actual solving is numerical.
\begin{center}
\begin{tabular}{|l|l|l|l|}
\hline
Model at $\Gamma$ & $S$ & $\beta$ & Band 2 \\
\hline
Only $\mathbf{G}=0$ (one wave) & $0$ & $1.25$ & $1342\,\mathrm{Hz}$ \\
Kirchhoff, all waves & $0.45$ & $1.14$ & $1280\,\mathrm{Hz}$ \\
RM, all waves & $0.50$ & $1.13$ & $1273.5\,\mathrm{Hz}$ \\
\hline
\end{tabular}
\end{center}
The plate bends under the resonator. That bending is a soft spring in series with
the resonator's spring, so the frequency drops. Of the Kirchhoff value $0.45$,
the first four waves contribute about $0.31$; each gives
$1/\bigl((4797/1280)^2 - 1\bigr) \approx 0.077$.

\paragraph{Step 2: look at $X$.}
At $X$ there are $d=2$ bare waves at $1194.4\,\mathrm{Hz}$,
$e^{+i\pi x/a}$ and $e^{-i\pi x/a}$. Two combinations:
\begin{itemize}
  \item \emph{sum}, proportional to $\cos(\pi x/a)$: antinode at $x_0 = 0$,
        couples to the resonator. Eq.~\eqref{eq:d_waves} with $d=2$ gives
        $846$ and $1694\,\mathrm{Hz}$; the full Eq.~\eqref{eq:pwe} gives
        $835.6$ and $1652.4\,\mathrm{Hz}$.
  \item \emph{difference}, proportional to $\sin(\pi x/a)$: node at $x_0 = 0$,
        so $w_0 = 0$ and the frequency stays bare, $1194.4\,\mathrm{Hz}$.
\end{itemize}
Sorted at $X$: $835.6$, $\mathbf{1194.4}$, $1652.4\,\mathrm{Hz}$. Band 2 at $X$ is
the sin pattern, and $1194.4 < 1273.5\,\mathrm{Hz}$, so $X$ is the lowest point
of band 2. \textbf{The upper edge is a pure Bragg frequency}: it contains neither
$m$ nor $k_s$.

% ---------------------------------------------------------------------
\section{Lower edge: from 1200 down to 1031.4\,Hz}

Band 1 peaks at $M$, where $d = 4$ and $\lambda_p = (2378.2\,\mathrm{Hz})^2$.
\begin{center}
\begin{tabular}{|p{0.36\textwidth}|l|l|p{0.08\textwidth}|}
\hline
Step & Equation & Band 1 at $M$ & Plate share \\
\hline
No lattice & \eqref{eq:one_wave}, $\lambda_p\to\infty$ & $1200\,\mathrm{Hz}$ & \\
Zone ends at $M$ (fixes A) & \eqref{eq:d_waves}, $d=1$ & $1153\,\mathrm{Hz}$ & \\
All four waves (fixes part of B) & \eqref{eq:d_waves}, $d=4$ & $1046\,\mathrm{Hz}$ & $5\,\%$ \\
All other waves too, RM (fixes B) & \eqref{eq:pwe} & $1031.4\,\mathrm{Hz}$ & \\
\hline
\end{tabular}
\end{center}
In the last step, Eq.~\eqref{eq:pwe} at $M$ reads
\begin{equation}
  1 = \mu_\mathrm{eff}\left(\frac{4\lambda}{\lambda_p - \lambda} + S_M\right) ,
\end{equation}
where $S_M$ is the sum over all waves outside the group of four. They all lie far
above $1031\,\mathrm{Hz}$, so $S_M > 0$ (spring-like). Back-calculated from the
solution, $S_M \approx 0.12$. The effect is the same as at $\Gamma$: extra
softness, lower frequency.

With only $5\,\%$ of the motion in the plate (Eq.~\eqref{eq:plate_share} with
$d=4$, $\beta = 0.760$), this state is almost pure resonator. That is why band 1
is so flat near $M$ in Fig.~\ref{fig:dispersion_vamm_1200}.

% ---------------------------------------------------------------------
\section{What kind of gap is this?}

Since the whole gap lies below $f_r$, we have $\beta < 1$ and
$m_\mathrm{eff} > 0$. The negative-mass region above $f_r$ plays no role.

\paragraph{Rigid-mass comparison.}
With a rigid mass, $m_\mathrm{eff} = m$, so Eq.~\eqref{eq:pwe} for $d$ waves gives
$1 = d\mu\lambda/(\lambda_p - \lambda)$. Multiplying out,
\begin{equation}
  \lambda_p - \lambda = d\mu\lambda
  \quad\Longrightarrow\quad
  \lambda = \frac{\lambda_p}{1 + d\mu} .
  \label{eq:rigid}
\end{equation}
For a complete gap, band 1 at $M$ must fall below band 2 at $X$:
\begin{equation}
  \frac{2378.2}{\sqrt{1 + 4\mu}} < 1194.4
  \;\Longleftrightarrow\; \sqrt{1+4\mu} > 1.991
  \;\Longleftrightarrow\; 1 + 4\mu > 3.964
  \;\Longleftrightarrow\; \mu > 0.74 .
\end{equation}
A rigid mass would need three quarters of the plate mass. The real $\mu$ is only
$0.25$, but the plate sees $\mu_\mathrm{eff}$:
\begin{center}
\begin{tabular}{|l|l|l|}
\hline
$f$ & $\beta$ & $\mu_\mathrm{eff} = 0.25/(1-\beta)$ \\
\hline
$1031.4\,\mathrm{Hz}$ & $0.739$ & $0.96$ \\
$1046\,\mathrm{Hz}$ & $0.760$ & $1.04$ \\
$1100\,\mathrm{Hz}$ & $0.840$ & $1.57$ \\
$1190\,\mathrm{Hz}$ & $0.983$ & $15$ \\
\hline
\end{tabular}
\end{center}
All are above $0.74$. Inserting the value $1.04$ at $1046\,\mathrm{Hz}$ into
Eq.~\eqref{eq:rigid} gives
\begin{equation}
  f = \frac{2378.2}{\sqrt{1 + 4\cdot 1.04}} = 1046\,\mathrm{Hz} .
\end{equation}
The $d=4$ line of the lower-edge table \emph{is} the rigid-mass formula, with the
mass evaluated at the frequency where it acts.

\paragraph{Result.}
This is a \emph{Bragg gap with a resonance-boosted mass}:
\begin{itemize}
  \item the \textbf{upper edge} comes from the lattice (sin pattern at $X$,
        untouched by the resonator);
  \item the \textbf{lower edge} comes from the resonator's growing effective
        mass pulling band 1 at $M$ down.
\end{itemize}
It works this well because $f_r$ sits just above the $X$ fold. The width is
$164\,\mathrm{Hz}$ ($14.7\,\%$, FE), compared with $142\,\mathrm{Hz}$ ($11\,\%$)
from the one-wave estimate.

% ---------------------------------------------------------------------
\section{When is the gap above or below $f_r$?}
\label{sec:gap_regimes}

A sweep of $f_r$ with the RM plane-wave model, $\mu = 0.25$, $12$ $\mathbf{k}$-points
per path segment, gap between bands 1 and 2. Edges are good to a few Hz (this
coarser sweep gives $1032\,\mathrm{Hz}$ where the converged value is
$1031.4\,\mathrm{Hz}$).
\begin{center}
\begin{tabular}{|l|l|l|l|l|}
\hline
$f_r$ [Hz] & Gap [Hz] & Width [Hz] & Gap relative to $f_r$ & Upper edge set by \\
\hline
500 & 487--554 & 67 & contains $f_r$ & resonator band \\
800 & 748--874 & 126 & contains $f_r$ & resonator band \\
1000 & 901--1076 & 176 & contains $f_r$ & resonator band \\
1100 & 969--1170 & 201 & contains $f_r$ & resonator band \\
1150 & 1001--1194 & 193 & contains $f_r$ & sin pattern at $X$ \\
1200 & 1032--1194 & 162 & below $f_r$ & sin pattern at $X$ \\
1300 & 1090--1194 & 105 & below $f_r$ & sin pattern at $X$ \\
1400 & 1142--1194 & 52 & below $f_r$ & sin pattern at $X$ \\
1500 & 1190--1194 & 5 & below $f_r$ & sin pattern at $X$ \\
1600 & none & -- & -- & -- \\
\hline
\end{tabular}
\end{center}

\paragraph{Rule 1: this gap can never lie entirely above $f_r$.}
Band 1 always stays below $f_r$. Take $\lambda$ between $0$ and the smaller of
$\lambda_r$ and the lowest bare $\lambda_n^0$ at this $\mathbf{k}$. Then
$\mu_\mathrm{eff} > 0$ and every term of the sum in Eq.~\eqref{eq:pwe} is
positive. At $\lambda = 0$ the right side is $0$. As $\lambda$ approaches the
smaller limit, the right side goes to $+\infty$, either through
$\mu_\mathrm{eff}$ or through the lowest term. In between it rises steadily, so
it equals $1$ exactly once. That root is band 1: below $f_r$ \emph{and} below the
bare plate band, at every $\mathbf{k}$.

\paragraph{Rule 2: the upper edge is the lower of two candidates.}
\begin{itemize}
  \item the resonator band near $\Gamma$, at about
        $f_r\sqrt{1+\mu-\mu S}$, which moves with $f_r$ and lies above $f_r$ as
        long as $S < 1$ (up to $f_r \approx 1765\,\mathrm{Hz}$ here);
  \item the sin pattern at $X$, fixed at $f_X$ and independent of the resonator.
\end{itemize}
The switch happens at roughly $f_r \approx 1120\,\mathrm{Hz}$. (In the sweep the
lowest point of the resonator band lies slightly off $\Gamma$ on $\Gamma$--$X$.)

\paragraph{The three regimes.}
\begin{enumerate}
  \item \textbf{$f_r$ below $f_X$: the gap contains $f_r$.} Band 1 flattens below
        $f_r$ and band 2 starts above it. This is the classical local-resonance
        gap once $f_r$ is well below the switch.
  \item \textbf{$f_r$ above $f_X$, up to about $1510\,\mathrm{Hz}$ here: the
        whole gap is below $f_r$.} The upper edge is stuck at $f_X$, and the lower
        edge is band 1 at $M$, pulled down by the growing $\mu_\mathrm{eff}$. The
        production case $f_r = 1200\,\mathrm{Hz}$ is just past the boundary.
  \item \textbf{$f_r$ above about $1510\,\mathrm{Hz}$: no complete gap.} Band 1
        at $M$ has risen past $f_X$.
\end{enumerate}
A gap above $f_r$ could only be a higher one, between bands 2 and 3. That gap did
not open for any tuning tested ($300$--$2000\,\mathrm{Hz}$), because band 2
reaches $2378\,\mathrm{Hz}$ at $M$. The widest gap found was at
$f_r = 1100\,\mathrm{Hz}$ ($201\,\mathrm{Hz}$), close to where the upper edge
switches candidates (candidate for the stage~7 parameter study).

% ---------------------------------------------------------------------
\section{Check questions with answers}

\paragraph{Question 1.} \emph{Double the resonator mass while keeping $f_r$
($\mu = 0.5$). What happens to each gap edge, and why does one of them not move?}

\emph{Answer.} The gap widens. The larger mass drags the lower edge down
(mass loading through $\mu_\mathrm{eff}$), from $1032$ to $927\,\mathrm{Hz}$. The
upper edge stays at $1194.4\,\mathrm{Hz}$. The reason is not that it is
sinusoidal (the cos pattern is too), but that the sin pattern has a \emph{node at
the resonator}: $w_0 = 0$, so no force acts on it. Width: $162 \to 268\,\mathrm{Hz}$.

\paragraph{Question 2.} \emph{Tune the resonator to $f_r = 800\,\mathrm{Hz}$, well
below the $X$ frequency of $1194\,\mathrm{Hz}$. Which mechanisms still act, what
now determines the edges, and does the gap sit below or above $f_r$?}

\emph{Answer.} This is different from the $1200\,\mathrm{Hz}$ case. The model
gives $748$--$874\,\mathrm{Hz}$:
\begin{itemize}
  \item it is a local-resonance gap, not mostly Bragg: neither edge is a Bragg
        frequency, and the sin pattern at $X$ lies far above the gap and plays no
        role;
  \item the gap \emph{straddles} $f_r$, starting below $800\,\mathrm{Hz}$ and
        ending above it;
  \item $\mu$ still matters, and now it moves \emph{both} edges: with
        $\mu = 0.5$ the gap becomes $707$--$936\,\mathrm{Hz}$. The upper edge is
        no longer pinned.
\end{itemize}

\paragraph{Question 3.} \emph{Why can the upper edge not rise above about
$1194\,\mathrm{Hz}$, however heavy or stiff the resonator is? What would you have
to change in the cell to raise it?}

\emph{Answer.} The upper-edge mode does not see the resonator, because it has a
node there. One option is a configuration where this mode does see a resonator.
Moving the single resonator does not achieve that: the two waves at $X$ always
have one combination that vanishes at $\mathbf{x}_0$, wherever $\mathbf{x}_0$ is.
One would need a \emph{second resonator per cell} at a different position, where
that pattern is not zero. The other route is to raise the bare frequency itself:
for bending waves $f_X \propto h/a^2$, so a smaller cell or a thicker plate.

% =====================================================================

\chapter{Why the gap sits at 1031--1195\,Hz and not at 1200--1342\,Hz}
\label{ch:gap_position}

The effective-medium estimate of Chapter~\ref{ch:local_resonance} predicts a
gap of $1200$--$1342\,\mathrm{Hz}$. The FE diagram shows a complete gap of
$1031$--$1195\,\mathrm{Hz}$. This chapter explains the difference. Every
equation is derived here, starting from Newton's law for the resonator, so the
chapter can be read on its own.

Unless stated otherwise, numbers come from the Reissner--Mindlin (RM)
plane-wave model, which agrees with FE to within $0.33\,\%$. The exact bare
frequency at $X$ is $1194.4\,\mathrm{Hz}$; earlier chapters round it to
$1190\,\mathrm{Hz}$. No conclusion there changes.

\begin{center}
\begin{tabular}{|l|p{0.62\textwidth}|}
\hline
Symbol & Meaning \\
\hline
$\omega$, $f$ & angular frequency and frequency of a wave (one dot of the diagram) \\
$\lambda = \omega^2$ & eigenvalue belonging to that dot \\
$\mathbf{k}$ & Bloch wavevector, the horizontal axis of the diagram \\
$\mathbf{x}$, $\mathbf{x}_0$ & position on the plate; position of the resonator (origin placed there, $\mathbf{x}_0 = 0$) \\
$a$, $A = a^2$ & cell size ($0.1\,\mathrm{m}$) and cell area \\
$\rho h$ & plate mass per area \\
$m$, $k_s$ & resonator mass and spring stiffness \\
$\lambda_r = \omega_r^2 = k_s/m$ & the resonator's own eigenvalue ($f_r = 1200\,\mathrm{Hz}$ unless stated) \\
$\mu = m/(\rho h A)$ & resonator mass over plate mass per cell, $0.25$ \\
$\beta = \lambda/\lambda_r = (f/f_r)^2$ & squared frequency relative to the resonator \\
$\varphi_n(\mathbf{x})$ & bare plate wave $n$, defined below \\
$\lambda_n^0$ & bare eigenvalue of wave $n$ (plate without resonator) \\
$\lambda_p$ & $\lambda_n^0$ of the one wave, or group of waves, kept in a simplified model \\
$q_n$ & amplitude of wave $n$ in the plate motion \\
$w_0$ & plate deflection at the resonator \\
$u$ & displacement of the resonator mass \\
$F$ & force the resonator exerts on the plate \\
$Z$, $C$ & force per deflection of the resonator; deflection per force of the plate \\
$S$ & sum over all waves except a chosen group, defined where it is used \\
$d$ & number of bare waves sharing the same $\lambda_n^0$ at a given $\mathbf{k}$ \\
$f_X$ & bare plate frequency at $X$, $1194.4\,\mathrm{Hz}$ \\
\hline
\end{tabular}
\end{center}

% ---------------------------------------------------------------------
\section{Recap: what the diagram contains}

A bare plate wave in the cell, for a given $\mathbf{k}$, is
\begin{equation}
  \varphi_n(\mathbf{x}) = \frac{e^{i(\mathbf{k}+\mathbf{G}_n)\cdot\mathbf{x}}}{\sqrt{\rho h A}} .
\end{equation}
The $\mathbf{G}_n = (2\pi/a)\,(\text{integer}, \text{integer})$ are the reciprocal
lattice vectors. Several waves fit at the same $\mathbf{k}$ because adding
$\mathbf{G}$ changes nothing from cell to cell. Each wave has its own bare
eigenvalue $\lambda_n^0$, set by its wavenumber $|\mathbf{k}+\mathbf{G}_n|$.

The factor $1/\sqrt{\rho h A}$ is chosen so that
$\int \rho h\,|\varphi_n|^2\,\mathrm{d}A = 1$ over the cell. Since
$|e^{i\dots}| = 1$,
\begin{equation}
  |\varphi_n(\mathbf{x}_0)|^2 = \frac{1}{\rho h A} \quad\text{for every wave } n.
  \label{eq:phi_squared}
\end{equation}

\begin{center}
\begin{tabular}{|l|l|l|l|}
\hline
Point & $\mathbf{k}$ & Bare waves there ($d$ of them) & Bare frequency \\
\hline
$\Gamma$ & $(0,0)$ & 1 wave, $\mathbf{G} = 0$ & $0\,\mathrm{Hz}$ \\
$\Gamma$ & $(0,0)$ & 4 waves, wavelength $a$ & $4715\,\mathrm{Hz}$ \\
$X$ & $(\pi/a, 0)$ & 2 waves, $(\pm\pi/a, 0)$ & $1194.4\,\mathrm{Hz}$ \\
$M$ & $(\pi/a, \pi/a)$ & 4 waves, $(\pm\pi/a, \pm\pi/a)$ & $2378.2\,\mathrm{Hz}$ \\
\hline
\end{tabular}
\end{center}

A \emph{complete gap} is a frequency interval with no dot at any $\mathbf{k}$.
It runs from the highest point of band 1 to the lowest point of band 2.

% ---------------------------------------------------------------------
\section{The resonator as a mass that depends on frequency}

Newton's law for the resonator mass: the spring force on it is
$k_s(w_0 - u)$, and for harmonic motion $\ddot{u} = -\omega^2 u$, so
\begin{align}
  -m\omega^2 u &= k_s(w_0 - u) \\
  u\,(k_s - m\omega^2) &= k_s w_0 \\
  u &= w_0\,\frac{k_s}{k_s - m\omega^2} .
\end{align}
Divide numerator and denominator by $k_s$. Since
$m\omega^2/k_s = \omega^2/\omega_r^2 = \beta$,
\begin{equation}
  u = \frac{w_0}{1 - \beta} .
  \label{eq:u_over_w}
\end{equation}
The force on the plate is equal and opposite to the force on the mass (Newton's
third law). The force on the mass is $m\ddot{u} = -m\omega^2 u$, so
\begin{equation}
  F = m\omega^2 u = \frac{m\,\omega^2 w_0}{1 - \beta} .
\end{equation}
A \emph{rigid} mass $M_\mathrm{rigid}$ glued to the plate moves with $w_0$, and
by the same reasoning exerts $F = M_\mathrm{rigid}\,\omega^2 w_0$. The two forces
have the same form. To the plate, the resonator therefore looks like a rigid
mass of size
\begin{equation}
  \boxed{\;m_\mathrm{eff} = \frac{m}{1-\beta}\;}
  \qquad\text{and, divided by } \rho h A, \qquad
  \boxed{\;\mu_\mathrm{eff} = \frac{\mu}{1-\beta}\;}
\end{equation}
Since $1 - \beta = (\lambda_r - \lambda)/\lambda_r$, an equivalent form is
\begin{equation}
  \mu_\mathrm{eff} = \frac{\mu\,\lambda_r}{\lambda_r - \lambda} .
  \label{eq:mu_eff_lambda}
\end{equation}
\begin{itemize}
  \item $f$ far below $f_r$ ($\beta \to 0$): $m_\mathrm{eff} \to m$; the spring
        is effectively rigid.
  \item $f$ rising towards $f_r$ ($\beta \to 1$): $m_\mathrm{eff}$ becomes very
        large and positive.
  \item $f$ above $f_r$ ($\beta > 1$): $m_\mathrm{eff}$ is negative.
\end{itemize}
In short, with $\lambda = \omega^2$,
\begin{equation}
  F = m_\mathrm{eff}\,\lambda\, w_0 .
  \label{eq:force}
\end{equation}

% ---------------------------------------------------------------------
\section{The plate driven by that force}

The plate with a point force $F$ at $\mathbf{x}_0$ obeys, with $\mathcal{L}$ the
plate stiffness operator,
\begin{equation}
  \mathcal{L} w - \rho h\,\lambda\, w = F\,\delta(\mathbf{x} - \mathbf{x}_0) .
\end{equation}
Expand the plate motion in bare waves, $w(\mathbf{x}) = \sum_j q_j\,\varphi_j(\mathbf{x})$,
multiply by $\varphi_n^*$ and integrate over the cell:
\begin{itemize}
  \item the stiffness term gives $\lambda_n^0\, q_n$ ($\varphi_n$ is a bare wave
        and the waves are orthogonal);
  \item the inertia term gives $\lambda\, q_n$ (normalization);
  \item on the right, the delta function picks the value at $\mathbf{x}_0$,
        giving $F\,\varphi_n(\mathbf{x}_0)^*$.
\end{itemize}
Therefore, for every wave $n$,
\begin{equation}
  (\lambda_n^0 - \lambda)\, q_n = F\,\varphi_n(\mathbf{x}_0)^* .
  \label{eq:modal}
\end{equation}
This says how strongly each wave responds to the point force.

\emph{Remark for Reissner--Mindlin:} there $\varphi_n$ also has rotation
components, and Eq.~\eqref{eq:modal} holds with the deflection part of
$\varphi_n$. The structure is identical, and the RM numbers include this.

% ---------------------------------------------------------------------
\section{Closing the loop: the plane-wave expansion condition}

Eq.~\eqref{eq:modal} gives the amplitude of each wave,
\begin{equation}
  q_n = \frac{F\,\varphi_n(\mathbf{x}_0)^*}{\lambda_n^0 - \lambda} ,
  \label{eq:qn}
\end{equation}
and the deflection at the resonator is the sum of all waves there,
\begin{equation}
  w_0 = \sum_n \varphi_n(\mathbf{x}_0)\, q_n .
  \label{eq:w0_sum}
\end{equation}

\paragraph{Insert \eqref{eq:qn} into \eqref{eq:w0_sum}.}
\begin{equation}
  w_0 = \sum_n \varphi_n(\mathbf{x}_0)\cdot\frac{F\,\varphi_n(\mathbf{x}_0)^*}{\lambda_n^0 - \lambda} .
\end{equation}
$F$ does not depend on $n$, so it comes out of the sum, and
$\varphi_n\varphi_n^* = |\varphi_n|^2$:
\begin{equation}
  w_0 = F\,\sum_n \frac{|\varphi_n(\mathbf{x}_0)|^2}{\lambda_n^0 - \lambda} .
\end{equation}
Define the sum as the \emph{compliance} $C$, the deflection per unit force:
\begin{equation}
  C = \sum_n \frac{|\varphi_n(\mathbf{x}_0)|^2}{\lambda_n^0 - \lambda} ,
  \qquad
  w_0 = C\,F .
  \label{eq:compliance}
\end{equation}
Likewise call the factor in Eq.~\eqref{eq:force} $Z = m_\mathrm{eff}\,\lambda$,
the force per unit deflection:
\begin{equation}
  F = Z\, w_0 .
  \label{eq:Z}
\end{equation}

\paragraph{Insert \eqref{eq:Z} into \eqref{eq:compliance} and divide by $w_0$.}
\begin{equation}
  w_0 = C\, Z\, w_0
  \quad\Longrightarrow\quad
  1 = Z\, C \qquad (w_0 \neq 0).
  \label{eq:ZC}
\end{equation}
The deflection creates a force, and that force creates exactly the same
deflection: the condition for a wave that sustains itself.

\paragraph{Rewriting $1 = ZC$ with mass ratios.}
Insert the definitions of $Z$ and $C$:
\begin{equation}
  1 = m_\mathrm{eff}\,\lambda\,\sum_n \frac{|\varphi_n(\mathbf{x}_0)|^2}{\lambda_n^0 - \lambda} .
\end{equation}
By Eq.~\eqref{eq:phi_squared}, $|\varphi_n(\mathbf{x}_0)|^2 = 1/(\rho h A)$ for
every $n$, so it comes out of the sum:
\begin{equation}
  1 = m_\mathrm{eff}\,\lambda\,\frac{1}{\rho h A}\,\sum_n \frac{1}{\lambda_n^0 - \lambda} .
\end{equation}
Group $m_\mathrm{eff}/(\rho h A) = \mu_\mathrm{eff}$ and move $\lambda$ into the
sum:
\begin{equation}
  \boxed{\;1 = \mu_\mathrm{eff}\,\sum_n \frac{\lambda}{\lambda_n^0 - \lambda}\;}
  \label{eq:pwe}
\end{equation}
This is the \emph{plane-wave expansion (PWE) condition}. At a given $\mathbf{k}$,
every $\lambda$ that satisfies it is one dot of the diagram. ``Kirchhoff'' or
``RM'' only decides which values $\lambda_n^0$ go into the sum; a ``Kirchhoff
compliance'' is $C$ evaluated with Kirchhoff values. In the notes the same
condition is written $1 = Z/A\cdot\sum_G g$, with
$g = 1/(\rho h(\omega_n^2 - \omega^2))$.

\paragraph{What if $w_0 = 0$?}
Then $F = 0$ by Eq.~\eqref{eq:Z}, and Eq.~\eqref{eq:modal} becomes
$(\lambda_n^0 - \lambda)\, q_n = 0$. A nonzero wave requires
$\lambda = \lambda_n^0$. A mode with a node at the resonator keeps its bare
frequency. Such modes do not appear in Eq.~\eqref{eq:pwe}, because we divided
by $w_0$, and must be added by hand.

% ---------------------------------------------------------------------
\section{Why a wave acts like a spring or like a mass}

Take one wave and write $g = F\,\varphi_n(\mathbf{x}_0)^*$ for the force it
feels. Eq.~\eqref{eq:qn} reads
\begin{equation}
  q_n = \frac{g}{\lambda_n^0 - \lambda} .
\end{equation}
Compare with the two elementary elements (unit mass, deflection $s$):
\begin{itemize}
  \item a \emph{spring} of stiffness $\kappa_s$: $\kappa_s\, s = g$, so
        $s = g/\kappa_s$. The deflection is along the force and does not depend
        on frequency.
  \item a \emph{mass} of $1$: $-\lambda\, s = g$, so $s = -g/\lambda$. The
        deflection is against the force.
\end{itemize}
Now the denominator decides:
\begin{itemize}
  \item \textbf{Wave above the driving frequency,} $\lambda_n^0 \gg \lambda$:
        $\lambda_n^0 - \lambda \approx \lambda_n^0 > 0$, so
        $q_n \approx g/\lambda_n^0$. This is exactly the spring formula with
        $\kappa_s = \lambda_n^0$.
  \item \textbf{Wave below the driving frequency,} $\lambda_n^0 \ll \lambda$:
        $\lambda_n^0 - \lambda \approx -\lambda < 0$, so
        $q_n \approx -g/\lambda$. This is exactly the mass formula.
\end{itemize}
Physically, each wave is an oscillator with a stiffness part $\lambda_n^0$ and
an inertia part $-\lambda$; the larger one decides its behaviour. For
$\lambda_n^0 = 0$ (no stiffness at all) the formula is exactly a mass.

% ---------------------------------------------------------------------
\section{One wave: where 1200--1342\,Hz comes from}

Keep only one term in Eq.~\eqref{eq:pwe}, with bare eigenvalue $\lambda_p$, and
use Eq.~\eqref{eq:mu_eff_lambda}:
\begin{equation}
  1 = \frac{\mu\,\lambda_r}{\lambda_r - \lambda}\cdot\frac{\lambda}{\lambda_p - \lambda} .
\end{equation}
Multiply both sides by $(\lambda_r - \lambda)(\lambda_p - \lambda)$:
\begin{equation}
  \boxed{\;(\lambda_p - \lambda)(\lambda_r - \lambda) = \mu\,\lambda_r\,\lambda\;}
  \label{eq:one_wave}
\end{equation}

\paragraph{Same equation as the determinant of Chapter~\ref{ch:local_resonance}.}
The two-level model with matrix
$H = \bigl(\begin{smallmatrix}\lambda_p+\mu\lambda_r & -\sqrt{\mu}\lambda_r\\ -\sqrt{\mu}\lambda_r & \lambda_r\end{smallmatrix}\bigr)$
gave $(\lambda_p + \mu\lambda_r - \lambda)(\lambda_r - \lambda) - \mu\lambda_r^2 = 0$.
Group the first bracket as $(\lambda_p - \lambda) + \mu\lambda_r$:
\begin{align}
  \bigl[(\lambda_p - \lambda) + \mu\lambda_r\bigr](\lambda_r - \lambda) - \mu\lambda_r^2 &= 0 \\
  (\lambda_p - \lambda)(\lambda_r - \lambda) + \mu\lambda_r(\lambda_r - \lambda) - \mu\lambda_r^2 &= 0 \\
  (\lambda_p - \lambda)(\lambda_r - \lambda) + \mu\lambda_r^2 - \mu\lambda_r\lambda - \mu\lambda_r^2 &= 0 \\
  (\lambda_p - \lambda)(\lambda_r - \lambda) - \mu\lambda_r\lambda &= 0 ,
\end{align}
which is Eq.~\eqref{eq:one_wave}. This form is useful because the left side
contains only the two uncoupled systems (it vanishes at $\lambda = \lambda_p$
and $\lambda = \lambda_r$), and the right side contains the coupling. With
$\mu = 0$ the uncoupled solutions come back.

\paragraph{The two bands.}
For each $\lambda_p$, Eq.~\eqref{eq:one_wave} has two solutions for $\lambda$.
Multiplying out gives a quadratic,
\begin{equation}
  \lambda^2 - \bigl(\lambda_p + \lambda_r(1+\mu)\bigr)\lambda + \lambda_p\lambda_r = 0 .
  \label{eq:quadratic}
\end{equation}
Call its two roots $\lambda_1 < \lambda_2$; they are the two bands at this $\mathbf{k}$. Writing the quadratic as $(\lambda - \lambda_1)(\lambda - \lambda_2) = \lambda^2 - (\lambda_1 + \lambda_2)\lambda + \lambda_1\lambda_2 = 0$ and comparing with Eq.~\eqref{eq:quadratic} gives $\lambda_1 + \lambda_2 = \lambda_p + \lambda_r(1+\mu)$ and $\lambda_1\lambda_2 = \lambda_p\lambda_r$. These two expressions are not the roots themselves, only their sum and their product.

\paragraph{Why two limits?}
The estimate pretends the plate is a uniform medium without a lattice. Then
$\lambda_p$ simply grows with the wavenumber, from $0$ at $\mathbf{k}=0$ towards
infinity. The bands as functions of $\lambda_p$, from Eq.~\eqref{eq:quadratic}:
\begin{center}
\begin{tabular}{|l|l|l|}
\hline
Bare plate frequency & Lower band & Upper band \\
\hline
$0\,\mathrm{Hz}$ & $0\,\mathrm{Hz}$ & $1342\,\mathrm{Hz}$ \\
$1000\,\mathrm{Hz}$ & $824\,\mathrm{Hz}$ & $1456\,\mathrm{Hz}$ \\
$2378\,\mathrm{Hz}$ & $1153\,\mathrm{Hz}$ & $2475\,\mathrm{Hz}$ \\
$10000\,\mathrm{Hz}$ & $1198\,\mathrm{Hz}$ & $10018\,\mathrm{Hz}$ \\
\hline
\end{tabular}
\end{center}
Both bands rise as $\lambda_p$ rises. The lower band is therefore \emph{highest}
at $\lambda_p \to \infty$ and the upper band \emph{lowest} at $\lambda_p = 0$.
These two ends bound the gap.

\paragraph{Limit $\lambda_p = 0$.}
Insert into Eq.~\eqref{eq:one_wave}:
$(-\lambda)(\lambda_r - \lambda) = \mu\lambda_r\lambda$. One solution is
$\lambda = 0$. Otherwise divide by $\lambda$:
\begin{align}
  -(\lambda_r - \lambda) &= \mu\lambda_r \\
  \lambda &= \lambda_r(1+\mu) \\
  f &= f_r\sqrt{1+\mu} = 1200\cdot\sqrt{1.25} = 1342\,\mathrm{Hz}.
\end{align}

\paragraph{Limit $\lambda_p \to \infty$.}
For very large $\lambda_p$ the sum of the roots is dominated by $\lambda_p$, so
the upper root is $\approx \lambda_p$. Then
\begin{equation}
  \text{lower root} = \frac{\text{product}}{\text{upper root}}
  \approx \frac{\lambda_p\lambda_r}{\lambda_p} = \lambda_r ,
\end{equation}
that is, $1200\,\mathrm{Hz}$.

\paragraph{The two hidden assumptions.}
\begin{itemize}
  \item \textbf{Assumption A:} $\lambda_p$ can grow without limit. In the
        lattice the zone ends at $M$, where the bare value stops at
        $2378\,\mathrm{Hz}$.
  \item \textbf{Assumption B:} one wave per $\mathbf{k}$ is enough. A point
        force excites all waves.
\end{itemize}

% ---------------------------------------------------------------------
\section{Several waves with the same bare frequency}

If $d$ waves share the same $\lambda_p$ and the others are ignored, the $d$ terms
in Eq.~\eqref{eq:pwe} are identical:
\begin{equation}
  1 = d\,\mu_\mathrm{eff}\,\frac{\lambda}{\lambda_p - \lambda} ,
\end{equation}
and the same multiplication gives
\begin{equation}
  \boxed{\;(\lambda_p - \lambda)(\lambda_r - \lambda) = d\,\mu\,\lambda_r\,\lambda\;}
  \label{eq:d_waves}
\end{equation}
This is Eq.~\eqref{eq:one_wave} with $\mu$ replaced by $d\mu$: $d$ waves act like
$d$ times the mass ratio. Of those $d$ waves, one combination has $w_0 \neq 0$;
the other $d-1$ combinations have a node at the resonator and stay at the bare
frequency.

\paragraph{Plate share of the motion.}
Let the coupling combination have amplitude $q$, normalized like the bare waves.
With $d$ waves adding at $\mathbf{x}_0$, $|w_0|^2 = q^2\,d/(\rho h A)$. The
kinetic energy in the plate is proportional to $q^2$, in the resonator to
$m\,u^2$. Using Eq.~\eqref{eq:u_over_w},
\begin{equation}
  m\,u^2 = \frac{m\, q^2\, d/(\rho h A)}{(1-\beta)^2} = q^2\,\frac{d\,\mu}{(1-\beta)^2} .
\end{equation}
Therefore
\begin{equation}
  \text{plate share} = \frac{q^2}{q^2 + m\,u^2} = \frac{1}{1 + d\mu/(1-\beta)^2} .
  \label{eq:plate_share}
\end{equation}
It is the same result as the eigenvector ratio of
Chapter~\ref{ch:local_resonance}, written in terms of $\beta$.

% ---------------------------------------------------------------------
\section{Upper edge: 1342 $\to$ 1273.5 $\to$ 1194.4\,Hz}

\paragraph{Step 1: stay at $\Gamma$, keep all waves.}
The sum in Eq.~\eqref{eq:pwe} has two parts:
\begin{itemize}
  \item the $\mathbf{G}=0$ wave has $\lambda_n^0 = 0$, and its term is
        $\lambda/(0-\lambda) = -1$: the mass-like part, the whole cell moving
        rigidly;
  \item all other waves lie far above ($4715\,\mathrm{Hz}$ and more). Their sum
        is given the name
        \begin{equation}
          S = \sum_{n \neq 0} \frac{\lambda}{\lambda_n^0 - \lambda} > 0 ,
        \end{equation}
        which is spring-like.
\end{itemize}
Eq.~\eqref{eq:pwe} becomes $1 = \mu_\mathrm{eff}(-1 + S)$. Insert
$\mu_\mathrm{eff} = \mu/(1-\beta)$ and multiply by $(1-\beta)$:
\begin{equation}
  1 - \beta = -\mu + \mu S
  \quad\Longrightarrow\quad
  \boxed{\;\beta = 1 + \mu - \mu S\;}
\end{equation}
$S$ itself depends on $\beta$, so the actual solving is numerical.
\begin{center}
\begin{tabular}{|l|l|l|l|}
\hline
Model at $\Gamma$ & $S$ & $\beta$ & Band 2 \\
\hline
Only $\mathbf{G}=0$ (one wave) & $0$ & $1.25$ & $1342\,\mathrm{Hz}$ \\
Kirchhoff, all waves & $0.45$ & $1.14$ & $1280\,\mathrm{Hz}$ \\
RM, all waves & $0.50$ & $1.13$ & $1273.5\,\mathrm{Hz}$ \\
\hline
\end{tabular}
\end{center}
The plate bends under the resonator. That bending is a soft spring in series with
the resonator's spring, so the frequency drops. Of the Kirchhoff value $0.45$,
the first four waves contribute about $0.31$; each gives
$1/\bigl((4797/1280)^2 - 1\bigr) \approx 0.077$.

\paragraph{Step 2: look at $X$.}
At $X$ there are $d=2$ bare waves at $1194.4\,\mathrm{Hz}$,
$e^{+i\pi x/a}$ and $e^{-i\pi x/a}$. Two combinations:
\begin{itemize}
  \item \emph{sum}, proportional to $\cos(\pi x/a)$: antinode at $x_0 = 0$,
        couples to the resonator. Eq.~\eqref{eq:d_waves} with $d=2$ gives
        $846$ and $1694\,\mathrm{Hz}$; the full Eq.~\eqref{eq:pwe} gives
        $835.6$ and $1652.4\,\mathrm{Hz}$.
  \item \emph{difference}, proportional to $\sin(\pi x/a)$: node at $x_0 = 0$,
        so $w_0 = 0$ and the frequency stays bare, $1194.4\,\mathrm{Hz}$.
\end{itemize}
Sorted at $X$: $835.6$, $\mathbf{1194.4}$, $1652.4\,\mathrm{Hz}$. Band 2 at $X$ is
the sin pattern, and $1194.4 < 1273.5\,\mathrm{Hz}$, so $X$ is the lowest point
of band 2. \textbf{The upper edge is a pure Bragg frequency}: it contains neither
$m$ nor $k_s$.

% ---------------------------------------------------------------------
\section{Lower edge: from 1200 down to 1031.4\,Hz}

Band 1 peaks at $M$, where $d = 4$ and $\lambda_p = (2378.2\,\mathrm{Hz})^2$.
\begin{center}
\begin{tabular}{|p{0.36\textwidth}|l|l|p{0.08\textwidth}|}
\hline
Step & Equation & Band 1 at $M$ & Plate share \\
\hline
No lattice & \eqref{eq:one_wave}, $\lambda_p\to\infty$ & $1200\,\mathrm{Hz}$ & \\
Zone ends at $M$ (fixes A) & \eqref{eq:d_waves}, $d=1$ & $1153\,\mathrm{Hz}$ & \\
All four waves (fixes part of B) & \eqref{eq:d_waves}, $d=4$ & $1046\,\mathrm{Hz}$ & $5\,\%$ \\
All other waves too, RM (fixes B) & \eqref{eq:pwe} & $1031.4\,\mathrm{Hz}$ & \\
\hline
\end{tabular}
\end{center}
In the last step, Eq.~\eqref{eq:pwe} at $M$ reads
\begin{equation}
  1 = \mu_\mathrm{eff}\left(\frac{4\lambda}{\lambda_p - \lambda} + S_M\right) ,
\end{equation}
where $S_M$ is the sum over all waves outside the group of four. They all lie far
above $1031\,\mathrm{Hz}$, so $S_M > 0$ (spring-like). Back-calculated from the
solution, $S_M \approx 0.12$. The effect is the same as at $\Gamma$: extra
softness, lower frequency.

With only $5\,\%$ of the motion in the plate (Eq.~\eqref{eq:plate_share} with
$d=4$, $\beta = 0.760$), this state is almost pure resonator. That is why band 1
is so flat near $M$ in Fig.~\ref{fig:dispersion_vamm_1200}.

% ---------------------------------------------------------------------
\section{What kind of gap is this?}

Since the whole gap lies below $f_r$, we have $\beta < 1$ and
$m_\mathrm{eff} > 0$. The negative-mass region above $f_r$ plays no role.

\paragraph{Rigid-mass comparison.}
With a rigid mass, $m_\mathrm{eff} = m$, so Eq.~\eqref{eq:pwe} for $d$ waves gives
$1 = d\mu\lambda/(\lambda_p - \lambda)$. Multiplying out,
\begin{equation}
  \lambda_p - \lambda = d\mu\lambda
  \quad\Longrightarrow\quad
  \lambda = \frac{\lambda_p}{1 + d\mu} .
  \label{eq:rigid}
\end{equation}
For a complete gap, band 1 at $M$ must fall below band 2 at $X$:
\begin{equation}
  \frac{2378.2}{\sqrt{1 + 4\mu}} < 1194.4
  \;\Longleftrightarrow\; \sqrt{1+4\mu} > 1.991
  \;\Longleftrightarrow\; 1 + 4\mu > 3.964
  \;\Longleftrightarrow\; \mu > 0.74 .
\end{equation}
A rigid mass would need three quarters of the plate mass. The real $\mu$ is only
$0.25$, but the plate sees $\mu_\mathrm{eff}$:
\begin{center}
\begin{tabular}{|l|l|l|}
\hline
$f$ & $\beta$ & $\mu_\mathrm{eff} = 0.25/(1-\beta)$ \\
\hline
$1031.4\,\mathrm{Hz}$ & $0.739$ & $0.96$ \\
$1046\,\mathrm{Hz}$ & $0.760$ & $1.04$ \\
$1100\,\mathrm{Hz}$ & $0.840$ & $1.57$ \\
$1190\,\mathrm{Hz}$ & $0.983$ & $15$ \\
\hline
\end{tabular}
\end{center}
All are above $0.74$. Inserting the value $1.04$ at $1046\,\mathrm{Hz}$ into
Eq.~\eqref{eq:rigid} gives
\begin{equation}
  f = \frac{2378.2}{\sqrt{1 + 4\cdot 1.04}} = 1046\,\mathrm{Hz} .
\end{equation}
The $d=4$ line of the lower-edge table \emph{is} the rigid-mass formula, with the
mass evaluated at the frequency where it acts.

\paragraph{Result.}
This is a \emph{Bragg gap with a resonance-boosted mass}:
\begin{itemize}
  \item the \textbf{upper edge} comes from the lattice (sin pattern at $X$,
        untouched by the resonator);
  \item the \textbf{lower edge} comes from the resonator's growing effective
        mass pulling band 1 at $M$ down.
\end{itemize}
It works this well because $f_r$ sits just above the $X$ fold. The width is
$164\,\mathrm{Hz}$ ($14.7\,\%$, FE), compared with $142\,\mathrm{Hz}$ ($11\,\%$)
from the one-wave estimate.

% ---------------------------------------------------------------------
\section{When is the gap above or below $f_r$?}
\label{sec:gap_regimes}

A sweep of $f_r$ with the RM plane-wave model, $\mu = 0.25$, $12$ $\mathbf{k}$-points
per path segment, gap between bands 1 and 2. Edges are good to a few Hz (this
coarser sweep gives $1032\,\mathrm{Hz}$ where the converged value is
$1031.4\,\mathrm{Hz}$).
\begin{center}
\begin{tabular}{|l|l|l|l|l|}
\hline
$f_r$ [Hz] & Gap [Hz] & Width [Hz] & Gap relative to $f_r$ & Upper edge set by \\
\hline
500 & 487--554 & 67 & contains $f_r$ & resonator band \\
800 & 748--874 & 126 & contains $f_r$ & resonator band \\
1000 & 901--1076 & 176 & contains $f_r$ & resonator band \\
1100 & 969--1170 & 201 & contains $f_r$ & resonator band \\
1150 & 1001--1194 & 193 & contains $f_r$ & sin pattern at $X$ \\
1200 & 1032--1194 & 162 & below $f_r$ & sin pattern at $X$ \\
1300 & 1090--1194 & 105 & below $f_r$ & sin pattern at $X$ \\
1400 & 1142--1194 & 52 & below $f_r$ & sin pattern at $X$ \\
1500 & 1190--1194 & 5 & below $f_r$ & sin pattern at $X$ \\
1600 & none & -- & -- & -- \\
\hline
\end{tabular}
\end{center}

\paragraph{Rule 1: this gap can never lie entirely above $f_r$.}
Band 1 always stays below $f_r$. Take $\lambda$ between $0$ and the smaller of
$\lambda_r$ and the lowest bare $\lambda_n^0$ at this $\mathbf{k}$. Then
$\mu_\mathrm{eff} > 0$ and every term of the sum in Eq.~\eqref{eq:pwe} is
positive. At $\lambda = 0$ the right side is $0$. As $\lambda$ approaches the
smaller limit, the right side goes to $+\infty$, either through
$\mu_\mathrm{eff}$ or through the lowest term. In between it rises steadily, so
it equals $1$ exactly once. That root is band 1: below $f_r$ \emph{and} below the
bare plate band, at every $\mathbf{k}$.

\paragraph{Rule 2: the upper edge is the lower of two candidates.}
\begin{itemize}
  \item the resonator band near $\Gamma$, at about
        $f_r\sqrt{1+\mu-\mu S}$, which moves with $f_r$ and lies above $f_r$ as
        long as $S < 1$ (up to $f_r \approx 1765\,\mathrm{Hz}$ here);
  \item the sin pattern at $X$, fixed at $f_X$ and independent of the resonator.
\end{itemize}
The switch happens at roughly $f_r \approx 1120\,\mathrm{Hz}$. (In the sweep the
lowest point of the resonator band lies slightly off $\Gamma$ on $\Gamma$--$X$.)

\paragraph{The three regimes.}
\begin{enumerate}
  \item \textbf{$f_r$ below $f_X$: the gap contains $f_r$.} Band 1 flattens below
        $f_r$ and band 2 starts above it. This is the classical local-resonance
        gap once $f_r$ is well below the switch.
  \item \textbf{$f_r$ above $f_X$, up to about $1510\,\mathrm{Hz}$ here: the
        whole gap is below $f_r$.} The upper edge is stuck at $f_X$, and the lower
        edge is band 1 at $M$, pulled down by the growing $\mu_\mathrm{eff}$. The
        production case $f_r = 1200\,\mathrm{Hz}$ is just past the boundary.
  \item \textbf{$f_r$ above about $1510\,\mathrm{Hz}$: no complete gap.} Band 1
        at $M$ has risen past $f_X$.
\end{enumerate}
A gap above $f_r$ could only be a higher one, between bands 2 and 3. That gap did
not open for any tuning tested ($300$--$2000\,\mathrm{Hz}$), because band 2
reaches $2378\,\mathrm{Hz}$ at $M$. The widest gap found was at
$f_r = 1100\,\mathrm{Hz}$ ($201\,\mathrm{Hz}$), close to where the upper edge
switches candidates (candidate for the stage~7 parameter study).

% ---------------------------------------------------------------------
\section{Check questions with answers}

\paragraph{Question 1.} \emph{Double the resonator mass while keeping $f_r$
($\mu = 0.5$). What happens to each gap edge, and why does one of them not move?}

\emph{Answer.} The gap widens. The larger mass drags the lower edge down
(mass loading through $\mu_\mathrm{eff}$), from $1032$ to $927\,\mathrm{Hz}$. The
upper edge stays at $1194.4\,\mathrm{Hz}$. The reason is not that it is
sinusoidal (the cos pattern is too), but that the sin pattern has a \emph{node at
the resonator}: $w_0 = 0$, so no force acts on it. Width: $162 \to 268\,\mathrm{Hz}$.

\paragraph{Question 2.} \emph{Tune the resonator to $f_r = 800\,\mathrm{Hz}$, well
below the $X$ frequency of $1194\,\mathrm{Hz}$. Which mechanisms still act, what
now determines the edges, and does the gap sit below or above $f_r$?}

\emph{Answer.} This is different from the $1200\,\mathrm{Hz}$ case. The model
gives $748$--$874\,\mathrm{Hz}$:
\begin{itemize}
  \item it is a local-resonance gap, not mostly Bragg: neither edge is a Bragg
        frequency, and the sin pattern at $X$ lies far above the gap and plays no
        role;
  \item the gap \emph{straddles} $f_r$, starting below $800\,\mathrm{Hz}$ and
        ending above it;
  \item $\mu$ still matters, and now it moves \emph{both} edges: with
        $\mu = 0.5$ the gap becomes $707$--$936\,\mathrm{Hz}$. The upper edge is
        no longer pinned.
\end{itemize}

\paragraph{Question 3.} \emph{Why can the upper edge not rise above about
$1194\,\mathrm{Hz}$, however heavy or stiff the resonator is? What would you have
to change in the cell to raise it?}

\emph{Answer.} The upper-edge mode does not see the resonator, because it has a
node there. One option is a configuration where this mode does see a resonator.
Moving the single resonator does not achieve that: the two waves at $X$ always
have one combination that vanishes at $\mathbf{x}_0$, wherever $\mathbf{x}_0$ is.
One would need a \emph{second resonator per cell} at a different position, where
that pattern is not zero. The other route is to raise the bare frequency itself:
for bending waves $f_X \propto h/a^2$, so a smaller cell or a thicker plate.

% =====================================================================

\chapter{Why the gap sits at 1031--1195\,Hz and not at 1200--1342\,Hz}
\label{ch:gap_position}

The effective-medium estimate of Chapter~\ref{ch:local_resonance} predicts a
gap of $1200$--$1342\,\mathrm{Hz}$. The FE diagram shows a complete gap of
$1031$--$1195\,\mathrm{Hz}$. This chapter explains the difference. Every
equation is derived here, starting from Newton's law for the resonator, so the
chapter can be read on its own.

Unless stated otherwise, numbers come from the Reissner--Mindlin (RM)
plane-wave model, which agrees with FE to within $0.33\,\%$. The exact bare
frequency at $X$ is $1194.4\,\mathrm{Hz}$; earlier chapters round it to
$1190\,\mathrm{Hz}$. No conclusion there changes.

\begin{center}
\begin{tabular}{|l|p{0.62\textwidth}|}
\hline
Symbol & Meaning \\
\hline
$\omega$, $f$ & angular frequency and frequency of a wave (one dot of the diagram) \\
$\lambda = \omega^2$ & eigenvalue belonging to that dot \\
$\mathbf{k}$ & Bloch wavevector, the horizontal axis of the diagram \\
$\mathbf{x}$, $\mathbf{x}_0$ & position on the plate; position of the resonator (origin placed there, $\mathbf{x}_0 = 0$) \\
$a$, $A = a^2$ & cell size ($0.1\,\mathrm{m}$) and cell area \\
$\rho h$ & plate mass per area \\
$m$, $k_s$ & resonator mass and spring stiffness \\
$\lambda_r = \omega_r^2 = k_s/m$ & the resonator's own eigenvalue ($f_r = 1200\,\mathrm{Hz}$ unless stated) \\
$\mu = m/(\rho h A)$ & resonator mass over plate mass per cell, $0.25$ \\
$\beta = \lambda/\lambda_r = (f/f_r)^2$ & squared frequency relative to the resonator \\
$\varphi_n(\mathbf{x})$ & bare plate wave $n$, defined below \\
$\lambda_n^0$ & bare eigenvalue of wave $n$ (plate without resonator) \\
$\lambda_p$ & $\lambda_n^0$ of the one wave, or group of waves, kept in a simplified model \\
$q_n$ & amplitude of wave $n$ in the plate motion \\
$w_0$ & plate deflection at the resonator \\
$u$ & displacement of the resonator mass \\
$F$ & force the resonator exerts on the plate \\
$Z$, $C$ & force per deflection of the resonator; deflection per force of the plate \\
$S$ & sum over all waves except a chosen group, defined where it is used \\
$d$ & number of bare waves sharing the same $\lambda_n^0$ at a given $\mathbf{k}$ \\
$f_X$ & bare plate frequency at $X$, $1194.4\,\mathrm{Hz}$ \\
\hline
\end{tabular}
\end{center}

% ---------------------------------------------------------------------
\section{Recap: what the diagram contains}

A bare plate wave in the cell, for a given $\mathbf{k}$, is
\begin{equation}
  \varphi_n(\mathbf{x}) = \frac{e^{i(\mathbf{k}+\mathbf{G}_n)\cdot\mathbf{x}}}{\sqrt{\rho h A}} .
\end{equation}
The $\mathbf{G}_n = (2\pi/a)\,(\text{integer}, \text{integer})$ are the reciprocal
lattice vectors. Several waves fit at the same $\mathbf{k}$ because adding
$\mathbf{G}$ changes nothing from cell to cell. Each wave has its own bare
eigenvalue $\lambda_n^0$, set by its wavenumber $|\mathbf{k}+\mathbf{G}_n|$.

The factor $1/\sqrt{\rho h A}$ is chosen so that
$\int \rho h\,|\varphi_n|^2\,\mathrm{d}A = 1$ over the cell. Since
$|e^{i\dots}| = 1$,
\begin{equation}
  |\varphi_n(\mathbf{x}_0)|^2 = \frac{1}{\rho h A} \quad\text{for every wave } n.
  \label{eq:phi_squared}
\end{equation}

\begin{center}
\begin{tabular}{|l|l|l|l|}
\hline
Point & $\mathbf{k}$ & Bare waves there ($d$ of them) & Bare frequency \\
\hline
$\Gamma$ & $(0,0)$ & 1 wave, $\mathbf{G} = 0$ & $0\,\mathrm{Hz}$ \\
$\Gamma$ & $(0,0)$ & 4 waves, wavelength $a$ & $4715\,\mathrm{Hz}$ \\
$X$ & $(\pi/a, 0)$ & 2 waves, $(\pm\pi/a, 0)$ & $1194.4\,\mathrm{Hz}$ \\
$M$ & $(\pi/a, \pi/a)$ & 4 waves, $(\pm\pi/a, \pm\pi/a)$ & $2378.2\,\mathrm{Hz}$ \\
\hline
\end{tabular}
\end{center}

A \emph{complete gap} is a frequency interval with no dot at any $\mathbf{k}$.
It runs from the highest point of band 1 to the lowest point of band 2.

% ---------------------------------------------------------------------
\section{The resonator as a mass that depends on frequency}

Newton's law for the resonator mass: the spring force on it is
$k_s(w_0 - u)$, and for harmonic motion $\ddot{u} = -\omega^2 u$, so
\begin{align}
  -m\omega^2 u &= k_s(w_0 - u) \\
  u\,(k_s - m\omega^2) &= k_s w_0 \\
  u &= w_0\,\frac{k_s}{k_s - m\omega^2} .
\end{align}
Divide numerator and denominator by $k_s$. Since
$m\omega^2/k_s = \omega^2/\omega_r^2 = \beta$,
\begin{equation}
  u = \frac{w_0}{1 - \beta} .
  \label{eq:u_over_w}
\end{equation}
The force on the plate is equal and opposite to the force on the mass (Newton's
third law). The force on the mass is $m\ddot{u} = -m\omega^2 u$, so
\begin{equation}
  F = m\omega^2 u = \frac{m\,\omega^2 w_0}{1 - \beta} .
\end{equation}
A \emph{rigid} mass $M_\mathrm{rigid}$ glued to the plate moves with $w_0$, and
by the same reasoning exerts $F = M_\mathrm{rigid}\,\omega^2 w_0$. The two forces
have the same form. To the plate, the resonator therefore looks like a rigid
mass of size
\begin{equation}
  \boxed{\;m_\mathrm{eff} = \frac{m}{1-\beta}\;}
  \qquad\text{and, divided by } \rho h A, \qquad
  \boxed{\;\mu_\mathrm{eff} = \frac{\mu}{1-\beta}\;}
\end{equation}
Since $1 - \beta = (\lambda_r - \lambda)/\lambda_r$, an equivalent form is
\begin{equation}
  \mu_\mathrm{eff} = \frac{\mu\,\lambda_r}{\lambda_r - \lambda} .
  \label{eq:mu_eff_lambda}
\end{equation}
\begin{itemize}
  \item $f$ far below $f_r$ ($\beta \to 0$): $m_\mathrm{eff} \to m$; the spring
        is effectively rigid.
  \item $f$ rising towards $f_r$ ($\beta \to 1$): $m_\mathrm{eff}$ becomes very
        large and positive.
  \item $f$ above $f_r$ ($\beta > 1$): $m_\mathrm{eff}$ is negative.
\end{itemize}
In short, with $\lambda = \omega^2$,
\begin{equation}
  F = m_\mathrm{eff}\,\lambda\, w_0 .
  \label{eq:force}
\end{equation}

% ---------------------------------------------------------------------
\section{The plate driven by that force}

The plate with a point force $F$ at $\mathbf{x}_0$ obeys, with $\mathcal{L}$ the
plate stiffness operator,
\begin{equation}
  \mathcal{L} w - \rho h\,\lambda\, w = F\,\delta(\mathbf{x} - \mathbf{x}_0) .
\end{equation}
Expand the plate motion in bare waves, $w(\mathbf{x}) = \sum_j q_j\,\varphi_j(\mathbf{x})$,
multiply by $\varphi_n^*$ and integrate over the cell:
\begin{itemize}
  \item the stiffness term gives $\lambda_n^0\, q_n$ ($\varphi_n$ is a bare wave
        and the waves are orthogonal);
  \item the inertia term gives $\lambda\, q_n$ (normalization);
  \item on the right, the delta function picks the value at $\mathbf{x}_0$,
        giving $F\,\varphi_n(\mathbf{x}_0)^*$.
\end{itemize}
Therefore, for every wave $n$,
\begin{equation}
  (\lambda_n^0 - \lambda)\, q_n = F\,\varphi_n(\mathbf{x}_0)^* .
  \label{eq:modal}
\end{equation}
This says how strongly each wave responds to the point force.

\emph{Remark for Reissner--Mindlin:} there $\varphi_n$ also has rotation
components, and Eq.~\eqref{eq:modal} holds with the deflection part of
$\varphi_n$. The structure is identical, and the RM numbers include this.

% ---------------------------------------------------------------------
\section{Closing the loop: the plane-wave expansion condition}

Eq.~\eqref{eq:modal} gives the amplitude of each wave,
\begin{equation}
  q_n = \frac{F\,\varphi_n(\mathbf{x}_0)^*}{\lambda_n^0 - \lambda} ,
  \label{eq:qn}
\end{equation}
and the deflection at the resonator is the sum of all waves there,
\begin{equation}
  w_0 = \sum_n \varphi_n(\mathbf{x}_0)\, q_n .
  \label{eq:w0_sum}
\end{equation}

\paragraph{Insert \eqref{eq:qn} into \eqref{eq:w0_sum}.}
\begin{equation}
  w_0 = \sum_n \varphi_n(\mathbf{x}_0)\cdot\frac{F\,\varphi_n(\mathbf{x}_0)^*}{\lambda_n^0 - \lambda} .
\end{equation}
$F$ does not depend on $n$, so it comes out of the sum, and
$\varphi_n\varphi_n^* = |\varphi_n|^2$:
\begin{equation}
  w_0 = F\,\sum_n \frac{|\varphi_n(\mathbf{x}_0)|^2}{\lambda_n^0 - \lambda} .
\end{equation}
Define the sum as the \emph{compliance} $C$, the deflection per unit force:
\begin{equation}
  C = \sum_n \frac{|\varphi_n(\mathbf{x}_0)|^2}{\lambda_n^0 - \lambda} ,
  \qquad
  w_0 = C\,F .
  \label{eq:compliance}
\end{equation}
Likewise call the factor in Eq.~\eqref{eq:force} $Z = m_\mathrm{eff}\,\lambda$,
the force per unit deflection:
\begin{equation}
  F = Z\, w_0 .
  \label{eq:Z}
\end{equation}

\paragraph{Insert \eqref{eq:Z} into \eqref{eq:compliance} and divide by $w_0$.}
\begin{equation}
  w_0 = C\, Z\, w_0
  \quad\Longrightarrow\quad
  1 = Z\, C \qquad (w_0 \neq 0).
  \label{eq:ZC}
\end{equation}
The deflection creates a force, and that force creates exactly the same
deflection: the condition for a wave that sustains itself.

\paragraph{Rewriting $1 = ZC$ with mass ratios.}
Insert the definitions of $Z$ and $C$:
\begin{equation}
  1 = m_\mathrm{eff}\,\lambda\,\sum_n \frac{|\varphi_n(\mathbf{x}_0)|^2}{\lambda_n^0 - \lambda} .
\end{equation}
By Eq.~\eqref{eq:phi_squared}, $|\varphi_n(\mathbf{x}_0)|^2 = 1/(\rho h A)$ for
every $n$, so it comes out of the sum:
\begin{equation}
  1 = m_\mathrm{eff}\,\lambda\,\frac{1}{\rho h A}\,\sum_n \frac{1}{\lambda_n^0 - \lambda} .
\end{equation}
Group $m_\mathrm{eff}/(\rho h A) = \mu_\mathrm{eff}$ and move $\lambda$ into the
sum:
\begin{equation}
  \boxed{\;1 = \mu_\mathrm{eff}\,\sum_n \frac{\lambda}{\lambda_n^0 - \lambda}\;}
  \label{eq:pwe}
\end{equation}
This is the \emph{plane-wave expansion (PWE) condition}. At a given $\mathbf{k}$,
every $\lambda$ that satisfies it is one dot of the diagram. ``Kirchhoff'' or
``RM'' only decides which values $\lambda_n^0$ go into the sum; a ``Kirchhoff
compliance'' is $C$ evaluated with Kirchhoff values. In the notes the same
condition is written $1 = Z/A\cdot\sum_G g$, with
$g = 1/(\rho h(\omega_n^2 - \omega^2))$.

\paragraph{What if $w_0 = 0$?}
Then $F = 0$ by Eq.~\eqref{eq:Z}, and Eq.~\eqref{eq:modal} becomes
$(\lambda_n^0 - \lambda)\, q_n = 0$. A nonzero wave requires
$\lambda = \lambda_n^0$. A mode with a node at the resonator keeps its bare
frequency. Such modes do not appear in Eq.~\eqref{eq:pwe}, because we divided
by $w_0$, and must be added by hand.

% ---------------------------------------------------------------------
\section{Why a wave acts like a spring or like a mass}

Take one wave and write $g = F\,\varphi_n(\mathbf{x}_0)^*$ for the force it
feels. Eq.~\eqref{eq:qn} reads
\begin{equation}
  q_n = \frac{g}{\lambda_n^0 - \lambda} .
\end{equation}
Compare with the two elementary elements (unit mass, deflection $s$):
\begin{itemize}
  \item a \emph{spring} of stiffness $\kappa_s$: $\kappa_s\, s = g$, so
        $s = g/\kappa_s$. The deflection is along the force and does not depend
        on frequency.
  \item a \emph{mass} of $1$: $-\lambda\, s = g$, so $s = -g/\lambda$. The
        deflection is against the force.
\end{itemize}
Now the denominator decides:
\begin{itemize}
  \item \textbf{Wave above the driving frequency,} $\lambda_n^0 \gg \lambda$:
        $\lambda_n^0 - \lambda \approx \lambda_n^0 > 0$, so
        $q_n \approx g/\lambda_n^0$. This is exactly the spring formula with
        $\kappa_s = \lambda_n^0$.
  \item \textbf{Wave below the driving frequency,} $\lambda_n^0 \ll \lambda$:
        $\lambda_n^0 - \lambda \approx -\lambda < 0$, so
        $q_n \approx -g/\lambda$. This is exactly the mass formula.
\end{itemize}
Physically, each wave is an oscillator with a stiffness part $\lambda_n^0$ and
an inertia part $-\lambda$; the larger one decides its behaviour. For
$\lambda_n^0 = 0$ (no stiffness at all) the formula is exactly a mass.

% ---------------------------------------------------------------------
\section{One wave: where 1200--1342\,Hz comes from}

Keep only one term in Eq.~\eqref{eq:pwe}, with bare eigenvalue $\lambda_p$, and
use Eq.~\eqref{eq:mu_eff_lambda}:
\begin{equation}
  1 = \frac{\mu\,\lambda_r}{\lambda_r - \lambda}\cdot\frac{\lambda}{\lambda_p - \lambda} .
\end{equation}
Multiply both sides by $(\lambda_r - \lambda)(\lambda_p - \lambda)$:
\begin{equation}
  \boxed{\;(\lambda_p - \lambda)(\lambda_r - \lambda) = \mu\,\lambda_r\,\lambda\;}
  \label{eq:one_wave}
\end{equation}

\paragraph{Same equation as the determinant of Chapter~\ref{ch:local_resonance}.}
The two-level model with matrix
$H = \bigl(\begin{smallmatrix}\lambda_p+\mu\lambda_r & -\sqrt{\mu}\lambda_r\\ -\sqrt{\mu}\lambda_r & \lambda_r\end{smallmatrix}\bigr)$
gave $(\lambda_p + \mu\lambda_r - \lambda)(\lambda_r - \lambda) - \mu\lambda_r^2 = 0$.
Group the first bracket as $(\lambda_p - \lambda) + \mu\lambda_r$:
\begin{align}
  \bigl[(\lambda_p - \lambda) + \mu\lambda_r\bigr](\lambda_r - \lambda) - \mu\lambda_r^2 &= 0 \\
  (\lambda_p - \lambda)(\lambda_r - \lambda) + \mu\lambda_r(\lambda_r - \lambda) - \mu\lambda_r^2 &= 0 \\
  (\lambda_p - \lambda)(\lambda_r - \lambda) + \mu\lambda_r^2 - \mu\lambda_r\lambda - \mu\lambda_r^2 &= 0 \\
  (\lambda_p - \lambda)(\lambda_r - \lambda) - \mu\lambda_r\lambda &= 0 ,
\end{align}
which is Eq.~\eqref{eq:one_wave}. This form is useful because the left side
contains only the two uncoupled systems (it vanishes at $\lambda = \lambda_p$
and $\lambda = \lambda_r$), and the right side contains the coupling. With
$\mu = 0$ the uncoupled solutions come back.

\paragraph{The two bands.}
For each $\lambda_p$, Eq.~\eqref{eq:one_wave} has two solutions for $\lambda$.
Multiplying out gives a quadratic,
\begin{equation}
  \lambda^2 - \bigl(\lambda_p + \lambda_r(1+\mu)\bigr)\lambda + \lambda_p\lambda_r = 0 .
  \label{eq:quadratic}
\end{equation}
Call its two roots $\lambda_1 < \lambda_2$; they are the two bands at this $\mathbf{k}$. Writing the quadratic as $(\lambda - \lambda_1)(\lambda - \lambda_2) = \lambda^2 - (\lambda_1 + \lambda_2)\lambda + \lambda_1\lambda_2 = 0$ and comparing with Eq.~\eqref{eq:quadratic} gives $\lambda_1 + \lambda_2 = \lambda_p + \lambda_r(1+\mu)$ and $\lambda_1\lambda_2 = \lambda_p\lambda_r$. These two expressions are not the roots themselves, only their sum and their product.

\paragraph{Why two limits?}
The estimate pretends the plate is a uniform medium without a lattice. Then
$\lambda_p$ simply grows with the wavenumber, from $0$ at $\mathbf{k}=0$ towards
infinity. The bands as functions of $\lambda_p$, from Eq.~\eqref{eq:quadratic}:
\begin{center}
\begin{tabular}{|l|l|l|}
\hline
Bare plate frequency & Lower band & Upper band \\
\hline
$0\,\mathrm{Hz}$ & $0\,\mathrm{Hz}$ & $1342\,\mathrm{Hz}$ \\
$1000\,\mathrm{Hz}$ & $824\,\mathrm{Hz}$ & $1456\,\mathrm{Hz}$ \\
$2378\,\mathrm{Hz}$ & $1153\,\mathrm{Hz}$ & $2475\,\mathrm{Hz}$ \\
$10000\,\mathrm{Hz}$ & $1198\,\mathrm{Hz}$ & $10018\,\mathrm{Hz}$ \\
\hline
\end{tabular}
\end{center}
Both bands rise as $\lambda_p$ rises. The lower band is therefore \emph{highest}
at $\lambda_p \to \infty$ and the upper band \emph{lowest} at $\lambda_p = 0$.
These two ends bound the gap.

\paragraph{Limit $\lambda_p = 0$.}
Insert into Eq.~\eqref{eq:one_wave}:
$(-\lambda)(\lambda_r - \lambda) = \mu\lambda_r\lambda$. One solution is
$\lambda = 0$. Otherwise divide by $\lambda$:
\begin{align}
  -(\lambda_r - \lambda) &= \mu\lambda_r \\
  \lambda &= \lambda_r(1+\mu) \\
  f &= f_r\sqrt{1+\mu} = 1200\cdot\sqrt{1.25} = 1342\,\mathrm{Hz}.
\end{align}

\paragraph{Limit $\lambda_p \to \infty$.}
For very large $\lambda_p$ the sum of the roots is dominated by $\lambda_p$, so
the upper root is $\approx \lambda_p$. Then
\begin{equation}
  \text{lower root} = \frac{\text{product}}{\text{upper root}}
  \approx \frac{\lambda_p\lambda_r}{\lambda_p} = \lambda_r ,
\end{equation}
that is, $1200\,\mathrm{Hz}$.

\paragraph{The two hidden assumptions.}
\begin{itemize}
  \item \textbf{Assumption A:} $\lambda_p$ can grow without limit. In the
        lattice the zone ends at $M$, where the bare value stops at
        $2378\,\mathrm{Hz}$.
  \item \textbf{Assumption B:} one wave per $\mathbf{k}$ is enough. A point
        force excites all waves.
\end{itemize}

% ---------------------------------------------------------------------
\section{Several waves with the same bare frequency}

If $d$ waves share the same $\lambda_p$ and the others are ignored, the $d$ terms
in Eq.~\eqref{eq:pwe} are identical:
\begin{equation}
  1 = d\,\mu_\mathrm{eff}\,\frac{\lambda}{\lambda_p - \lambda} ,
\end{equation}
and the same multiplication gives
\begin{equation}
  \boxed{\;(\lambda_p - \lambda)(\lambda_r - \lambda) = d\,\mu\,\lambda_r\,\lambda\;}
  \label{eq:d_waves}
\end{equation}
This is Eq.~\eqref{eq:one_wave} with $\mu$ replaced by $d\mu$: $d$ waves act like
$d$ times the mass ratio. Of those $d$ waves, one combination has $w_0 \neq 0$;
the other $d-1$ combinations have a node at the resonator and stay at the bare
frequency.

\paragraph{Plate share of the motion.}
Let the coupling combination have amplitude $q$, normalized like the bare waves.
With $d$ waves adding at $\mathbf{x}_0$, $|w_0|^2 = q^2\,d/(\rho h A)$. The
kinetic energy in the plate is proportional to $q^2$, in the resonator to
$m\,u^2$. Using Eq.~\eqref{eq:u_over_w},
\begin{equation}
  m\,u^2 = \frac{m\, q^2\, d/(\rho h A)}{(1-\beta)^2} = q^2\,\frac{d\,\mu}{(1-\beta)^2} .
\end{equation}
Therefore
\begin{equation}
  \text{plate share} = \frac{q^2}{q^2 + m\,u^2} = \frac{1}{1 + d\mu/(1-\beta)^2} .
  \label{eq:plate_share}
\end{equation}
It is the same result as the eigenvector ratio of
Chapter~\ref{ch:local_resonance}, written in terms of $\beta$.

% ---------------------------------------------------------------------
\section{Upper edge: 1342 $\to$ 1273.5 $\to$ 1194.4\,Hz}

\paragraph{Step 1: stay at $\Gamma$, keep all waves.}
The sum in Eq.~\eqref{eq:pwe} has two parts:
\begin{itemize}
  \item the $\mathbf{G}=0$ wave has $\lambda_n^0 = 0$, and its term is
        $\lambda/(0-\lambda) = -1$: the mass-like part, the whole cell moving
        rigidly;
  \item all other waves lie far above ($4715\,\mathrm{Hz}$ and more). Their sum
        is given the name
        \begin{equation}
          S = \sum_{n \neq 0} \frac{\lambda}{\lambda_n^0 - \lambda} > 0 ,
        \end{equation}
        which is spring-like.
\end{itemize}
Eq.~\eqref{eq:pwe} becomes $1 = \mu_\mathrm{eff}(-1 + S)$. Insert
$\mu_\mathrm{eff} = \mu/(1-\beta)$ and multiply by $(1-\beta)$:
\begin{equation}
  1 - \beta = -\mu + \mu S
  \quad\Longrightarrow\quad
  \boxed{\;\beta = 1 + \mu - \mu S\;}
\end{equation}
$S$ itself depends on $\beta$, so the actual solving is numerical.
\begin{center}
\begin{tabular}{|l|l|l|l|}
\hline
Model at $\Gamma$ & $S$ & $\beta$ & Band 2 \\
\hline
Only $\mathbf{G}=0$ (one wave) & $0$ & $1.25$ & $1342\,\mathrm{Hz}$ \\
Kirchhoff, all waves & $0.45$ & $1.14$ & $1280\,\mathrm{Hz}$ \\
RM, all waves & $0.50$ & $1.13$ & $1273.5\,\mathrm{Hz}$ \\
\hline
\end{tabular}
\end{center}
The plate bends under the resonator. That bending is a soft spring in series with
the resonator's spring, so the frequency drops. Of the Kirchhoff value $0.45$,
the first four waves contribute about $0.31$; each gives
$1/\bigl((4797/1280)^2 - 1\bigr) \approx 0.077$.

\paragraph{Step 2: look at $X$.}
At $X$ there are $d=2$ bare waves at $1194.4\,\mathrm{Hz}$,
$e^{+i\pi x/a}$ and $e^{-i\pi x/a}$. Two combinations:
\begin{itemize}
  \item \emph{sum}, proportional to $\cos(\pi x/a)$: antinode at $x_0 = 0$,
        couples to the resonator. Eq.~\eqref{eq:d_waves} with $d=2$ gives
        $846$ and $1694\,\mathrm{Hz}$; the full Eq.~\eqref{eq:pwe} gives
        $835.6$ and $1652.4\,\mathrm{Hz}$.
  \item \emph{difference}, proportional to $\sin(\pi x/a)$: node at $x_0 = 0$,
        so $w_0 = 0$ and the frequency stays bare, $1194.4\,\mathrm{Hz}$.
\end{itemize}
Sorted at $X$: $835.6$, $\mathbf{1194.4}$, $1652.4\,\mathrm{Hz}$. Band 2 at $X$ is
the sin pattern, and $1194.4 < 1273.5\,\mathrm{Hz}$, so $X$ is the lowest point
of band 2. \textbf{The upper edge is a pure Bragg frequency}: it contains neither
$m$ nor $k_s$.

% ---------------------------------------------------------------------
\section{Lower edge: from 1200 down to 1031.4\,Hz}

Band 1 peaks at $M$, where $d = 4$ and $\lambda_p = (2378.2\,\mathrm{Hz})^2$.
\begin{center}
\begin{tabular}{|p{0.36\textwidth}|l|l|p{0.08\textwidth}|}
\hline
Step & Equation & Band 1 at $M$ & Plate share \\
\hline
No lattice & \eqref{eq:one_wave}, $\lambda_p\to\infty$ & $1200\,\mathrm{Hz}$ & \\
Zone ends at $M$ (fixes A) & \eqref{eq:d_waves}, $d=1$ & $1153\,\mathrm{Hz}$ & \\
All four waves (fixes part of B) & \eqref{eq:d_waves}, $d=4$ & $1046\,\mathrm{Hz}$ & $5\,\%$ \\
All other waves too, RM (fixes B) & \eqref{eq:pwe} & $1031.4\,\mathrm{Hz}$ & \\
\hline
\end{tabular}
\end{center}
In the last step, Eq.~\eqref{eq:pwe} at $M$ reads
\begin{equation}
  1 = \mu_\mathrm{eff}\left(\frac{4\lambda}{\lambda_p - \lambda} + S_M\right) ,
\end{equation}
where $S_M$ is the sum over all waves outside the group of four. They all lie far
above $1031\,\mathrm{Hz}$, so $S_M > 0$ (spring-like). Back-calculated from the
solution, $S_M \approx 0.12$. The effect is the same as at $\Gamma$: extra
softness, lower frequency.

With only $5\,\%$ of the motion in the plate (Eq.~\eqref{eq:plate_share} with
$d=4$, $\beta = 0.760$), this state is almost pure resonator. That is why band 1
is so flat near $M$ in Fig.~\ref{fig:dispersion_vamm_1200}.

% ---------------------------------------------------------------------
\section{What kind of gap is this?}

Since the whole gap lies below $f_r$, we have $\beta < 1$ and
$m_\mathrm{eff} > 0$. The negative-mass region above $f_r$ plays no role.

\paragraph{Rigid-mass comparison.}
With a rigid mass, $m_\mathrm{eff} = m$, so Eq.~\eqref{eq:pwe} for $d$ waves gives
$1 = d\mu\lambda/(\lambda_p - \lambda)$. Multiplying out,
\begin{equation}
  \lambda_p - \lambda = d\mu\lambda
  \quad\Longrightarrow\quad
  \lambda = \frac{\lambda_p}{1 + d\mu} .
  \label{eq:rigid}
\end{equation}
For a complete gap, band 1 at $M$ must fall below band 2 at $X$:
\begin{equation}
  \frac{2378.2}{\sqrt{1 + 4\mu}} < 1194.4
  \;\Longleftrightarrow\; \sqrt{1+4\mu} > 1.991
  \;\Longleftrightarrow\; 1 + 4\mu > 3.964
  \;\Longleftrightarrow\; \mu > 0.74 .
\end{equation}
A rigid mass would need three quarters of the plate mass. The real $\mu$ is only
$0.25$, but the plate sees $\mu_\mathrm{eff}$:
\begin{center}
\begin{tabular}{|l|l|l|}
\hline
$f$ & $\beta$ & $\mu_\mathrm{eff} = 0.25/(1-\beta)$ \\
\hline
$1031.4\,\mathrm{Hz}$ & $0.739$ & $0.96$ \\
$1046\,\mathrm{Hz}$ & $0.760$ & $1.04$ \\
$1100\,\mathrm{Hz}$ & $0.840$ & $1.57$ \\
$1190\,\mathrm{Hz}$ & $0.983$ & $15$ \\
\hline
\end{tabular}
\end{center}
All are above $0.74$. Inserting the value $1.04$ at $1046\,\mathrm{Hz}$ into
Eq.~\eqref{eq:rigid} gives
\begin{equation}
  f = \frac{2378.2}{\sqrt{1 + 4\cdot 1.04}} = 1046\,\mathrm{Hz} .
\end{equation}
The $d=4$ line of the lower-edge table \emph{is} the rigid-mass formula, with the
mass evaluated at the frequency where it acts.

\paragraph{Result.}
This is a \emph{Bragg gap with a resonance-boosted mass}:
\begin{itemize}
  \item the \textbf{upper edge} comes from the lattice (sin pattern at $X$,
        untouched by the resonator);
  \item the \textbf{lower edge} comes from the resonator's growing effective
        mass pulling band 1 at $M$ down.
\end{itemize}
It works this well because $f_r$ sits just above the $X$ fold. The width is
$164\,\mathrm{Hz}$ ($14.7\,\%$, FE), compared with $142\,\mathrm{Hz}$ ($11\,\%$)
from the one-wave estimate.

% ---------------------------------------------------------------------
\section{When is the gap above or below $f_r$?}
\label{sec:gap_regimes}

A sweep of $f_r$ with the RM plane-wave model, $\mu = 0.25$, $12$ $\mathbf{k}$-points
per path segment, gap between bands 1 and 2. Edges are good to a few Hz (this
coarser sweep gives $1032\,\mathrm{Hz}$ where the converged value is
$1031.4\,\mathrm{Hz}$).
\begin{center}
\begin{tabular}{|l|l|l|l|l|}
\hline
$f_r$ [Hz] & Gap [Hz] & Width [Hz] & Gap relative to $f_r$ & Upper edge set by \\
\hline
500 & 487--554 & 67 & contains $f_r$ & resonator band \\
800 & 748--874 & 126 & contains $f_r$ & resonator band \\
1000 & 901--1076 & 176 & contains $f_r$ & resonator band \\
1100 & 969--1170 & 201 & contains $f_r$ & resonator band \\
1150 & 1001--1194 & 193 & contains $f_r$ & sin pattern at $X$ \\
1200 & 1032--1194 & 162 & below $f_r$ & sin pattern at $X$ \\
1300 & 1090--1194 & 105 & below $f_r$ & sin pattern at $X$ \\
1400 & 1142--1194 & 52 & below $f_r$ & sin pattern at $X$ \\
1500 & 1190--1194 & 5 & below $f_r$ & sin pattern at $X$ \\
1600 & none & -- & -- & -- \\
\hline
\end{tabular}
\end{center}

\paragraph{Rule 1: this gap can never lie entirely above $f_r$.}
Band 1 always stays below $f_r$. Take $\lambda$ between $0$ and the smaller of
$\lambda_r$ and the lowest bare $\lambda_n^0$ at this $\mathbf{k}$. Then
$\mu_\mathrm{eff} > 0$ and every term of the sum in Eq.~\eqref{eq:pwe} is
positive. At $\lambda = 0$ the right side is $0$. As $\lambda$ approaches the
smaller limit, the right side goes to $+\infty$, either through
$\mu_\mathrm{eff}$ or through the lowest term. In between it rises steadily, so
it equals $1$ exactly once. That root is band 1: below $f_r$ \emph{and} below the
bare plate band, at every $\mathbf{k}$.

\paragraph{Rule 2: the upper edge is the lower of two candidates.}
\begin{itemize}
  \item the resonator band near $\Gamma$, at about
        $f_r\sqrt{1+\mu-\mu S}$, which moves with $f_r$ and lies above $f_r$ as
        long as $S < 1$ (up to $f_r \approx 1765\,\mathrm{Hz}$ here);
  \item the sin pattern at $X$, fixed at $f_X$ and independent of the resonator.
\end{itemize}
The switch happens at roughly $f_r \approx 1120\,\mathrm{Hz}$. (In the sweep the
lowest point of the resonator band lies slightly off $\Gamma$ on $\Gamma$--$X$.)

\paragraph{The three regimes.}
\begin{enumerate}
  \item \textbf{$f_r$ below $f_X$: the gap contains $f_r$.} Band 1 flattens below
        $f_r$ and band 2 starts above it. This is the classical local-resonance
        gap once $f_r$ is well below the switch.
  \item \textbf{$f_r$ above $f_X$, up to about $1510\,\mathrm{Hz}$ here: the
        whole gap is below $f_r$.} The upper edge is stuck at $f_X$, and the lower
        edge is band 1 at $M$, pulled down by the growing $\mu_\mathrm{eff}$. The
        production case $f_r = 1200\,\mathrm{Hz}$ is just past the boundary.
  \item \textbf{$f_r$ above about $1510\,\mathrm{Hz}$: no complete gap.} Band 1
        at $M$ has risen past $f_X$.
\end{enumerate}
A gap above $f_r$ could only be a higher one, between bands 2 and 3. That gap did
not open for any tuning tested ($300$--$2000\,\mathrm{Hz}$), because band 2
reaches $2378\,\mathrm{Hz}$ at $M$. The widest gap found was at
$f_r = 1100\,\mathrm{Hz}$ ($201\,\mathrm{Hz}$), close to where the upper edge
switches candidates (candidate for the stage~7 parameter study).

% ---------------------------------------------------------------------
\section{Check questions with answers}

\paragraph{Question 1.} \emph{Double the resonator mass while keeping $f_r$
($\mu = 0.5$). What happens to each gap edge, and why does one of them not move?}

\emph{Answer.} The gap widens. The larger mass drags the lower edge down
(mass loading through $\mu_\mathrm{eff}$), from $1032$ to $927\,\mathrm{Hz}$. The
upper edge stays at $1194.4\,\mathrm{Hz}$. The reason is not that it is
sinusoidal (the cos pattern is too), but that the sin pattern has a \emph{node at
the resonator}: $w_0 = 0$, so no force acts on it. Width: $162 \to 268\,\mathrm{Hz}$.

\paragraph{Question 2.} \emph{Tune the resonator to $f_r = 800\,\mathrm{Hz}$, well
below the $X$ frequency of $1194\,\mathrm{Hz}$. Which mechanisms still act, what
now determines the edges, and does the gap sit below or above $f_r$?}

\emph{Answer.} This is different from the $1200\,\mathrm{Hz}$ case. The model
gives $748$--$874\,\mathrm{Hz}$:
\begin{itemize}
  \item it is a local-resonance gap, not mostly Bragg: neither edge is a Bragg
        frequency, and the sin pattern at $X$ lies far above the gap and plays no
        role;
  \item the gap \emph{straddles} $f_r$, starting below $800\,\mathrm{Hz}$ and
        ending above it;
  \item $\mu$ still matters, and now it moves \emph{both} edges: with
        $\mu = 0.5$ the gap becomes $707$--$936\,\mathrm{Hz}$. The upper edge is
        no longer pinned.
\end{itemize}

\paragraph{Question 3.} \emph{Why can the upper edge not rise above about
$1194\,\mathrm{Hz}$, however heavy or stiff the resonator is? What would you have
to change in the cell to raise it?}

\emph{Answer.} The upper-edge mode does not see the resonator, because it has a
node there. One option is a configuration where this mode does see a resonator.
Moving the single resonator does not achieve that: the two waves at $X$ always
have one combination that vanishes at $\mathbf{x}_0$, wherever $\mathbf{x}_0$ is.
One would need a \emph{second resonator per cell} at a different position, where
that pattern is not zero. The other route is to raise the bare frequency itself:
for bending waves $f_X \propto h/a^2$, so a smaller cell or a thicker plate.

% =====================================================================

\chapter{Why the gap sits at 1031--1195\,Hz and not at 1200--1342\,Hz}
\label{ch:gap_position}

The effective-medium estimate of Chapter~\ref{ch:local_resonance} predicts a
gap of $1200$--$1342\,\mathrm{Hz}$. The FE diagram shows a complete gap of
$1031$--$1195\,\mathrm{Hz}$. This chapter explains the difference. Every
equation is derived here, starting from Newton's law for the resonator, so the
chapter can be read on its own.

Unless stated otherwise, numbers come from the Reissner--Mindlin (RM)
plane-wave model, which agrees with FE to within $0.33\,\%$. The exact bare
frequency at $X$ is $1194.4\,\mathrm{Hz}$; earlier chapters round it to
$1190\,\mathrm{Hz}$. No conclusion there changes.

\begin{center}
\begin{tabular}{|l|p{0.62\textwidth}|}
\hline
Symbol & Meaning \\
\hline
$\omega$, $f$ & angular frequency and frequency of a wave (one dot of the diagram) \\
$\lambda = \omega^2$ & eigenvalue belonging to that dot \\
$\mathbf{k}$ & Bloch wavevector, the horizontal axis of the diagram \\
$\mathbf{x}$, $\mathbf{x}_0$ & position on the plate; position of the resonator (origin placed there, $\mathbf{x}_0 = 0$) \\
$a$, $A = a^2$ & cell size ($0.1\,\mathrm{m}$) and cell area \\
$\rho h$ & plate mass per area \\
$m$, $k_s$ & resonator mass and spring stiffness \\
$\lambda_r = \omega_r^2 = k_s/m$ & the resonator's own eigenvalue ($f_r = 1200\,\mathrm{Hz}$ unless stated) \\
$\mu = m/(\rho h A)$ & resonator mass over plate mass per cell, $0.25$ \\
$\beta = \lambda/\lambda_r = (f/f_r)^2$ & squared frequency relative to the resonator \\
$\varphi_n(\mathbf{x})$ & bare plate wave $n$, defined below \\
$\lambda_n^0$ & bare eigenvalue of wave $n$ (plate without resonator) \\
$\lambda_p$ & $\lambda_n^0$ of the one wave, or group of waves, kept in a simplified model \\
$q_n$ & amplitude of wave $n$ in the plate motion \\
$w_0$ & plate deflection at the resonator \\
$u$ & displacement of the resonator mass \\
$F$ & force the resonator exerts on the plate \\
$Z$, $C$ & force per deflection of the resonator; deflection per force of the plate \\
$S$ & sum over all waves except a chosen group, defined where it is used \\
$d$ & number of bare waves sharing the same $\lambda_n^0$ at a given $\mathbf{k}$ \\
$f_X$ & bare plate frequency at $X$, $1194.4\,\mathrm{Hz}$ \\
\hline
\end{tabular}
\end{center}

% ---------------------------------------------------------------------
\section{Recap: what the diagram contains}

A bare plate wave in the cell, for a given $\mathbf{k}$, is
\begin{equation}
  \varphi_n(\mathbf{x}) = \frac{e^{i(\mathbf{k}+\mathbf{G}_n)\cdot\mathbf{x}}}{\sqrt{\rho h A}} .
\end{equation}
The $\mathbf{G}_n = (2\pi/a)\,(\text{integer}, \text{integer})$ are the reciprocal
lattice vectors. Several waves fit at the same $\mathbf{k}$ because adding
$\mathbf{G}$ changes nothing from cell to cell. Each wave has its own bare
eigenvalue $\lambda_n^0$, set by its wavenumber $|\mathbf{k}+\mathbf{G}_n|$.

The factor $1/\sqrt{\rho h A}$ is chosen so that
$\int \rho h\,|\varphi_n|^2\,\mathrm{d}A = 1$ over the cell. Since
$|e^{i\dots}| = 1$,
\begin{equation}
  |\varphi_n(\mathbf{x}_0)|^2 = \frac{1}{\rho h A} \quad\text{for every wave } n.
  \label{eq:phi_squared}
\end{equation}

\begin{center}
\begin{tabular}{|l|l|l|l|}
\hline
Point & $\mathbf{k}$ & Bare waves there ($d$ of them) & Bare frequency \\
\hline
$\Gamma$ & $(0,0)$ & 1 wave, $\mathbf{G} = 0$ & $0\,\mathrm{Hz}$ \\
$\Gamma$ & $(0,0)$ & 4 waves, wavelength $a$ & $4715\,\mathrm{Hz}$ \\
$X$ & $(\pi/a, 0)$ & 2 waves, $(\pm\pi/a, 0)$ & $1194.4\,\mathrm{Hz}$ \\
$M$ & $(\pi/a, \pi/a)$ & 4 waves, $(\pm\pi/a, \pm\pi/a)$ & $2378.2\,\mathrm{Hz}$ \\
\hline
\end{tabular}
\end{center}

A \emph{complete gap} is a frequency interval with no dot at any $\mathbf{k}$.
It runs from the highest point of band 1 to the lowest point of band 2.

% ---------------------------------------------------------------------
\section{The resonator as a mass that depends on frequency}

Newton's law for the resonator mass: the spring force on it is
$k_s(w_0 - u)$, and for harmonic motion $\ddot{u} = -\omega^2 u$, so
\begin{align}
  -m\omega^2 u &= k_s(w_0 - u) \\
  u\,(k_s - m\omega^2) &= k_s w_0 \\
  u &= w_0\,\frac{k_s}{k_s - m\omega^2} .
\end{align}
Divide numerator and denominator by $k_s$. Since
$m\omega^2/k_s = \omega^2/\omega_r^2 = \beta$,
\begin{equation}
  u = \frac{w_0}{1 - \beta} .
  \label{eq:u_over_w}
\end{equation}
The force on the plate is equal and opposite to the force on the mass (Newton's
third law). The force on the mass is $m\ddot{u} = -m\omega^2 u$, so
\begin{equation}
  F = m\omega^2 u = \frac{m\,\omega^2 w_0}{1 - \beta} .
\end{equation}
A \emph{rigid} mass $M_\mathrm{rigid}$ glued to the plate moves with $w_0$, and
by the same reasoning exerts $F = M_\mathrm{rigid}\,\omega^2 w_0$. The two forces
have the same form. To the plate, the resonator therefore looks like a rigid
mass of size
\begin{equation}
  \boxed{\;m_\mathrm{eff} = \frac{m}{1-\beta}\;}
  \qquad\text{and, divided by } \rho h A, \qquad
  \boxed{\;\mu_\mathrm{eff} = \frac{\mu}{1-\beta}\;}
\end{equation}
Since $1 - \beta = (\lambda_r - \lambda)/\lambda_r$, an equivalent form is
\begin{equation}
  \mu_\mathrm{eff} = \frac{\mu\,\lambda_r}{\lambda_r - \lambda} .
  \label{eq:mu_eff_lambda}
\end{equation}
\begin{itemize}
  \item $f$ far below $f_r$ ($\beta \to 0$): $m_\mathrm{eff} \to m$; the spring
        is effectively rigid.
  \item $f$ rising towards $f_r$ ($\beta \to 1$): $m_\mathrm{eff}$ becomes very
        large and positive.
  \item $f$ above $f_r$ ($\beta > 1$): $m_\mathrm{eff}$ is negative.
\end{itemize}
In short, with $\lambda = \omega^2$,
\begin{equation}
  F = m_\mathrm{eff}\,\lambda\, w_0 .
  \label{eq:force}
\end{equation}

% ---------------------------------------------------------------------
\section{The plate driven by that force}

The plate with a point force $F$ at $\mathbf{x}_0$ obeys, with $\mathcal{L}$ the
plate stiffness operator,
\begin{equation}
  \mathcal{L} w - \rho h\,\lambda\, w = F\,\delta(\mathbf{x} - \mathbf{x}_0) .
\end{equation}
Expand the plate motion in bare waves, $w(\mathbf{x}) = \sum_j q_j\,\varphi_j(\mathbf{x})$,
multiply by $\varphi_n^*$ and integrate over the cell:
\begin{itemize}
  \item the stiffness term gives $\lambda_n^0\, q_n$ ($\varphi_n$ is a bare wave
        and the waves are orthogonal);
  \item the inertia term gives $\lambda\, q_n$ (normalization);
  \item on the right, the delta function picks the value at $\mathbf{x}_0$,
        giving $F\,\varphi_n(\mathbf{x}_0)^*$.
\end{itemize}
Therefore, for every wave $n$,
\begin{equation}
  (\lambda_n^0 - \lambda)\, q_n = F\,\varphi_n(\mathbf{x}_0)^* .
  \label{eq:modal}
\end{equation}
This says how strongly each wave responds to the point force.

\emph{Remark for Reissner--Mindlin:} there $\varphi_n$ also has rotation
components, and Eq.~\eqref{eq:modal} holds with the deflection part of
$\varphi_n$. The structure is identical, and the RM numbers include this.

% ---------------------------------------------------------------------
\section{Closing the loop: the plane-wave expansion condition}

Eq.~\eqref{eq:modal} gives the amplitude of each wave,
\begin{equation}
  q_n = \frac{F\,\varphi_n(\mathbf{x}_0)^*}{\lambda_n^0 - \lambda} ,
  \label{eq:qn}
\end{equation}
and the deflection at the resonator is the sum of all waves there,
\begin{equation}
  w_0 = \sum_n \varphi_n(\mathbf{x}_0)\, q_n .
  \label{eq:w0_sum}
\end{equation}

\paragraph{Insert \eqref{eq:qn} into \eqref{eq:w0_sum}.}
\begin{equation}
  w_0 = \sum_n \varphi_n(\mathbf{x}_0)\cdot\frac{F\,\varphi_n(\mathbf{x}_0)^*}{\lambda_n^0 - \lambda} .
\end{equation}
$F$ does not depend on $n$, so it comes out of the sum, and
$\varphi_n\varphi_n^* = |\varphi_n|^2$:
\begin{equation}
  w_0 = F\,\sum_n \frac{|\varphi_n(\mathbf{x}_0)|^2}{\lambda_n^0 - \lambda} .
\end{equation}
Define the sum as the \emph{compliance} $C$, the deflection per unit force:
\begin{equation}
  C = \sum_n \frac{|\varphi_n(\mathbf{x}_0)|^2}{\lambda_n^0 - \lambda} ,
  \qquad
  w_0 = C\,F .
  \label{eq:compliance}
\end{equation}
Likewise call the factor in Eq.~\eqref{eq:force} $Z = m_\mathrm{eff}\,\lambda$,
the force per unit deflection:
\begin{equation}
  F = Z\, w_0 .
  \label{eq:Z}
\end{equation}

\paragraph{Insert \eqref{eq:Z} into \eqref{eq:compliance} and divide by $w_0$.}
\begin{equation}
  w_0 = C\, Z\, w_0
  \quad\Longrightarrow\quad
  1 = Z\, C \qquad (w_0 \neq 0).
  \label{eq:ZC}
\end{equation}
The deflection creates a force, and that force creates exactly the same
deflection: the condition for a wave that sustains itself.

\paragraph{Rewriting $1 = ZC$ with mass ratios.}
Insert the definitions of $Z$ and $C$:
\begin{equation}
  1 = m_\mathrm{eff}\,\lambda\,\sum_n \frac{|\varphi_n(\mathbf{x}_0)|^2}{\lambda_n^0 - \lambda} .
\end{equation}
By Eq.~\eqref{eq:phi_squared}, $|\varphi_n(\mathbf{x}_0)|^2 = 1/(\rho h A)$ for
every $n$, so it comes out of the sum:
\begin{equation}
  1 = m_\mathrm{eff}\,\lambda\,\frac{1}{\rho h A}\,\sum_n \frac{1}{\lambda_n^0 - \lambda} .
\end{equation}
Group $m_\mathrm{eff}/(\rho h A) = \mu_\mathrm{eff}$ and move $\lambda$ into the
sum:
\begin{equation}
  \boxed{\;1 = \mu_\mathrm{eff}\,\sum_n \frac{\lambda}{\lambda_n^0 - \lambda}\;}
  \label{eq:pwe}
\end{equation}
This is the \emph{plane-wave expansion (PWE) condition}. At a given $\mathbf{k}$,
every $\lambda$ that satisfies it is one dot of the diagram. ``Kirchhoff'' or
``RM'' only decides which values $\lambda_n^0$ go into the sum; a ``Kirchhoff
compliance'' is $C$ evaluated with Kirchhoff values. In the notes the same
condition is written $1 = Z/A\cdot\sum_G g$, with
$g = 1/(\rho h(\omega_n^2 - \omega^2))$.

\paragraph{What if $w_0 = 0$?}
Then $F = 0$ by Eq.~\eqref{eq:Z}, and Eq.~\eqref{eq:modal} becomes
$(\lambda_n^0 - \lambda)\, q_n = 0$. A nonzero wave requires
$\lambda = \lambda_n^0$. A mode with a node at the resonator keeps its bare
frequency. Such modes do not appear in Eq.~\eqref{eq:pwe}, because we divided
by $w_0$, and must be added by hand.

% ---------------------------------------------------------------------
\section{Why a wave acts like a spring or like a mass}

Take one wave and write $g = F\,\varphi_n(\mathbf{x}_0)^*$ for the force it
feels. Eq.~\eqref{eq:qn} reads
\begin{equation}
  q_n = \frac{g}{\lambda_n^0 - \lambda} .
\end{equation}
Compare with the two elementary elements (unit mass, deflection $s$):
\begin{itemize}
  \item a \emph{spring} of stiffness $\kappa_s$: $\kappa_s\, s = g$, so
        $s = g/\kappa_s$. The deflection is along the force and does not depend
        on frequency.
  \item a \emph{mass} of $1$: $-\lambda\, s = g$, so $s = -g/\lambda$. The
        deflection is against the force.
\end{itemize}
Now the denominator decides:
\begin{itemize}
  \item \textbf{Wave above the driving frequency,} $\lambda_n^0 \gg \lambda$:
        $\lambda_n^0 - \lambda \approx \lambda_n^0 > 0$, so
        $q_n \approx g/\lambda_n^0$. This is exactly the spring formula with
        $\kappa_s = \lambda_n^0$.
  \item \textbf{Wave below the driving frequency,} $\lambda_n^0 \ll \lambda$:
        $\lambda_n^0 - \lambda \approx -\lambda < 0$, so
        $q_n \approx -g/\lambda$. This is exactly the mass formula.
\end{itemize}
Physically, each wave is an oscillator with a stiffness part $\lambda_n^0$ and
an inertia part $-\lambda$; the larger one decides its behaviour. For
$\lambda_n^0 = 0$ (no stiffness at all) the formula is exactly a mass.

% ---------------------------------------------------------------------
\section{One wave: where 1200--1342\,Hz comes from}

Keep only one term in Eq.~\eqref{eq:pwe}, with bare eigenvalue $\lambda_p$, and
use Eq.~\eqref{eq:mu_eff_lambda}:
\begin{equation}
  1 = \frac{\mu\,\lambda_r}{\lambda_r - \lambda}\cdot\frac{\lambda}{\lambda_p - \lambda} .
\end{equation}
Multiply both sides by $(\lambda_r - \lambda)(\lambda_p - \lambda)$:
\begin{equation}
  \boxed{\;(\lambda_p - \lambda)(\lambda_r - \lambda) = \mu\,\lambda_r\,\lambda\;}
  \label{eq:one_wave}
\end{equation}

\paragraph{Same equation as the determinant of Chapter~\ref{ch:local_resonance}.}
The two-level model with matrix
$H = \bigl(\begin{smallmatrix}\lambda_p+\mu\lambda_r & -\sqrt{\mu}\lambda_r\\ -\sqrt{\mu}\lambda_r & \lambda_r\end{smallmatrix}\bigr)$
gave $(\lambda_p + \mu\lambda_r - \lambda)(\lambda_r - \lambda) - \mu\lambda_r^2 = 0$.
Group the first bracket as $(\lambda_p - \lambda) + \mu\lambda_r$:
\begin{align}
  \bigl[(\lambda_p - \lambda) + \mu\lambda_r\bigr](\lambda_r - \lambda) - \mu\lambda_r^2 &= 0 \\
  (\lambda_p - \lambda)(\lambda_r - \lambda) + \mu\lambda_r(\lambda_r - \lambda) - \mu\lambda_r^2 &= 0 \\
  (\lambda_p - \lambda)(\lambda_r - \lambda) + \mu\lambda_r^2 - \mu\lambda_r\lambda - \mu\lambda_r^2 &= 0 \\
  (\lambda_p - \lambda)(\lambda_r - \lambda) - \mu\lambda_r\lambda &= 0 ,
\end{align}
which is Eq.~\eqref{eq:one_wave}. This form is useful because the left side
contains only the two uncoupled systems (it vanishes at $\lambda = \lambda_p$
and $\lambda = \lambda_r$), and the right side contains the coupling. With
$\mu = 0$ the uncoupled solutions come back.

\paragraph{The two bands.}
For each $\lambda_p$, Eq.~\eqref{eq:one_wave} has two solutions for $\lambda$.
Multiplying out gives a quadratic,
\begin{equation}
  \lambda^2 - \bigl(\lambda_p + \lambda_r(1+\mu)\bigr)\lambda + \lambda_p\lambda_r = 0 .
  \label{eq:quadratic}
\end{equation}
Call its two roots $\lambda_1 < \lambda_2$; they are the two bands at this $\mathbf{k}$. Writing the quadratic as $(\lambda - \lambda_1)(\lambda - \lambda_2) = \lambda^2 - (\lambda_1 + \lambda_2)\lambda + \lambda_1\lambda_2 = 0$ and comparing with Eq.~\eqref{eq:quadratic} gives $\lambda_1 + \lambda_2 = \lambda_p + \lambda_r(1+\mu)$ and $\lambda_1\lambda_2 = \lambda_p\lambda_r$. These two expressions are not the roots themselves, only their sum and their product.

\paragraph{Why two limits?}
The estimate pretends the plate is a uniform medium without a lattice. Then
$\lambda_p$ simply grows with the wavenumber, from $0$ at $\mathbf{k}=0$ towards
infinity. The bands as functions of $\lambda_p$, from Eq.~\eqref{eq:quadratic}:
\begin{center}
\begin{tabular}{|l|l|l|}
\hline
Bare plate frequency & Lower band & Upper band \\
\hline
$0\,\mathrm{Hz}$ & $0\,\mathrm{Hz}$ & $1342\,\mathrm{Hz}$ \\
$1000\,\mathrm{Hz}$ & $824\,\mathrm{Hz}$ & $1456\,\mathrm{Hz}$ \\
$2378\,\mathrm{Hz}$ & $1153\,\mathrm{Hz}$ & $2475\,\mathrm{Hz}$ \\
$10000\,\mathrm{Hz}$ & $1198\,\mathrm{Hz}$ & $10018\,\mathrm{Hz}$ \\
\hline
\end{tabular}
\end{center}
Both bands rise as $\lambda_p$ rises. The lower band is therefore \emph{highest}
at $\lambda_p \to \infty$ and the upper band \emph{lowest} at $\lambda_p = 0$.
These two ends bound the gap.

\paragraph{Limit $\lambda_p = 0$.}
Insert into Eq.~\eqref{eq:one_wave}:
$(-\lambda)(\lambda_r - \lambda) = \mu\lambda_r\lambda$. One solution is
$\lambda = 0$. Otherwise divide by $\lambda$:
\begin{align}
  -(\lambda_r - \lambda) &= \mu\lambda_r \\
  \lambda &= \lambda_r(1+\mu) \\
  f &= f_r\sqrt{1+\mu} = 1200\cdot\sqrt{1.25} = 1342\,\mathrm{Hz}.
\end{align}

\paragraph{Limit $\lambda_p \to \infty$.}
For very large $\lambda_p$ the sum of the roots is dominated by $\lambda_p$, so
the upper root is $\approx \lambda_p$. Then
\begin{equation}
  \text{lower root} = \frac{\text{product}}{\text{upper root}}
  \approx \frac{\lambda_p\lambda_r}{\lambda_p} = \lambda_r ,
\end{equation}
that is, $1200\,\mathrm{Hz}$.

\paragraph{The two hidden assumptions.}
\begin{itemize}
  \item \textbf{Assumption A:} $\lambda_p$ can grow without limit. In the
        lattice the zone ends at $M$, where the bare value stops at
        $2378\,\mathrm{Hz}$.
  \item \textbf{Assumption B:} one wave per $\mathbf{k}$ is enough. A point
        force excites all waves.
\end{itemize}

% ---------------------------------------------------------------------
\section{Several waves with the same bare frequency}

If $d$ waves share the same $\lambda_p$ and the others are ignored, the $d$ terms
in Eq.~\eqref{eq:pwe} are identical:
\begin{equation}
  1 = d\,\mu_\mathrm{eff}\,\frac{\lambda}{\lambda_p - \lambda} ,
\end{equation}
and the same multiplication gives
\begin{equation}
  \boxed{\;(\lambda_p - \lambda)(\lambda_r - \lambda) = d\,\mu\,\lambda_r\,\lambda\;}
  \label{eq:d_waves}
\end{equation}
This is Eq.~\eqref{eq:one_wave} with $\mu$ replaced by $d\mu$: $d$ waves act like
$d$ times the mass ratio. Of those $d$ waves, one combination has $w_0 \neq 0$;
the other $d-1$ combinations have a node at the resonator and stay at the bare
frequency.

\paragraph{Plate share of the motion.}
Let the coupling combination have amplitude $q$, normalized like the bare waves.
With $d$ waves adding at $\mathbf{x}_0$, $|w_0|^2 = q^2\,d/(\rho h A)$. The
kinetic energy in the plate is proportional to $q^2$, in the resonator to
$m\,u^2$. Using Eq.~\eqref{eq:u_over_w},
\begin{equation}
  m\,u^2 = \frac{m\, q^2\, d/(\rho h A)}{(1-\beta)^2} = q^2\,\frac{d\,\mu}{(1-\beta)^2} .
\end{equation}
Therefore
\begin{equation}
  \text{plate share} = \frac{q^2}{q^2 + m\,u^2} = \frac{1}{1 + d\mu/(1-\beta)^2} .
  \label{eq:plate_share}
\end{equation}
It is the same result as the eigenvector ratio of
Chapter~\ref{ch:local_resonance}, written in terms of $\beta$.

% ---------------------------------------------------------------------
\section{Upper edge: 1342 $\to$ 1273.5 $\to$ 1194.4\,Hz}

\paragraph{Step 1: stay at $\Gamma$, keep all waves.}
The sum in Eq.~\eqref{eq:pwe} has two parts:
\begin{itemize}
  \item the $\mathbf{G}=0$ wave has $\lambda_n^0 = 0$, and its term is
        $\lambda/(0-\lambda) = -1$: the mass-like part, the whole cell moving
        rigidly;
  \item all other waves lie far above ($4715\,\mathrm{Hz}$ and more). Their sum
        is given the name
        \begin{equation}
          S = \sum_{n \neq 0} \frac{\lambda}{\lambda_n^0 - \lambda} > 0 ,
        \end{equation}
        which is spring-like.
\end{itemize}
Eq.~\eqref{eq:pwe} becomes $1 = \mu_\mathrm{eff}(-1 + S)$. Insert
$\mu_\mathrm{eff} = \mu/(1-\beta)$ and multiply by $(1-\beta)$:
\begin{equation}
  1 - \beta = -\mu + \mu S
  \quad\Longrightarrow\quad
  \boxed{\;\beta = 1 + \mu - \mu S\;}
\end{equation}
$S$ itself depends on $\beta$, so the actual solving is numerical.
\begin{center}
\begin{tabular}{|l|l|l|l|}
\hline
Model at $\Gamma$ & $S$ & $\beta$ & Band 2 \\
\hline
Only $\mathbf{G}=0$ (one wave) & $0$ & $1.25$ & $1342\,\mathrm{Hz}$ \\
Kirchhoff, all waves & $0.45$ & $1.14$ & $1280\,\mathrm{Hz}$ \\
RM, all waves & $0.50$ & $1.13$ & $1273.5\,\mathrm{Hz}$ \\
\hline
\end{tabular}
\end{center}
The plate bends under the resonator. That bending is a soft spring in series with
the resonator's spring, so the frequency drops. Of the Kirchhoff value $0.45$,
the first four waves contribute about $0.31$; each gives
$1/\bigl((4797/1280)^2 - 1\bigr) \approx 0.077$.

\paragraph{Step 2: look at $X$.}
At $X$ there are $d=2$ bare waves at $1194.4\,\mathrm{Hz}$,
$e^{+i\pi x/a}$ and $e^{-i\pi x/a}$. Two combinations:
\begin{itemize}
  \item \emph{sum}, proportional to $\cos(\pi x/a)$: antinode at $x_0 = 0$,
        couples to the resonator. Eq.~\eqref{eq:d_waves} with $d=2$ gives
        $846$ and $1694\,\mathrm{Hz}$; the full Eq.~\eqref{eq:pwe} gives
        $835.6$ and $1652.4\,\mathrm{Hz}$.
  \item \emph{difference}, proportional to $\sin(\pi x/a)$: node at $x_0 = 0$,
        so $w_0 = 0$ and the frequency stays bare, $1194.4\,\mathrm{Hz}$.
\end{itemize}
Sorted at $X$: $835.6$, $\mathbf{1194.4}$, $1652.4\,\mathrm{Hz}$. Band 2 at $X$ is
the sin pattern, and $1194.4 < 1273.5\,\mathrm{Hz}$, so $X$ is the lowest point
of band 2. \textbf{The upper edge is a pure Bragg frequency}: it contains neither
$m$ nor $k_s$.

% ---------------------------------------------------------------------
\section{Lower edge: from 1200 down to 1031.4\,Hz}

Band 1 peaks at $M$, where $d = 4$ and $\lambda_p = (2378.2\,\mathrm{Hz})^2$.
\begin{center}
\begin{tabular}{|p{0.36\textwidth}|l|l|p{0.08\textwidth}|}
\hline
Step & Equation & Band 1 at $M$ & Plate share \\
\hline
No lattice & \eqref{eq:one_wave}, $\lambda_p\to\infty$ & $1200\,\mathrm{Hz}$ & \\
Zone ends at $M$ (fixes A) & \eqref{eq:d_waves}, $d=1$ & $1153\,\mathrm{Hz}$ & \\
All four waves (fixes part of B) & \eqref{eq:d_waves}, $d=4$ & $1046\,\mathrm{Hz}$ & $5\,\%$ \\
All other waves too, RM (fixes B) & \eqref{eq:pwe} & $1031.4\,\mathrm{Hz}$ & \\
\hline
\end{tabular}
\end{center}
In the last step, Eq.~\eqref{eq:pwe} at $M$ reads
\begin{equation}
  1 = \mu_\mathrm{eff}\left(\frac{4\lambda}{\lambda_p - \lambda} + S_M\right) ,
\end{equation}
where $S_M$ is the sum over all waves outside the group of four. They all lie far
above $1031\,\mathrm{Hz}$, so $S_M > 0$ (spring-like). Back-calculated from the
solution, $S_M \approx 0.12$. The effect is the same as at $\Gamma$: extra
softness, lower frequency.

With only $5\,\%$ of the motion in the plate (Eq.~\eqref{eq:plate_share} with
$d=4$, $\beta = 0.760$), this state is almost pure resonator. That is why band 1
is so flat near $M$ in Fig.~\ref{fig:dispersion_vamm_1200}.

% ---------------------------------------------------------------------
\section{What kind of gap is this?}

Since the whole gap lies below $f_r$, we have $\beta < 1$ and
$m_\mathrm{eff} > 0$. The negative-mass region above $f_r$ plays no role.

\paragraph{Rigid-mass comparison.}
With a rigid mass, $m_\mathrm{eff} = m$, so Eq.~\eqref{eq:pwe} for $d$ waves gives
$1 = d\mu\lambda/(\lambda_p - \lambda)$. Multiplying out,
\begin{equation}
  \lambda_p - \lambda = d\mu\lambda
  \quad\Longrightarrow\quad
  \lambda = \frac{\lambda_p}{1 + d\mu} .
  \label{eq:rigid}
\end{equation}
For a complete gap, band 1 at $M$ must fall below band 2 at $X$:
\begin{equation}
  \frac{2378.2}{\sqrt{1 + 4\mu}} < 1194.4
  \;\Longleftrightarrow\; \sqrt{1+4\mu} > 1.991
  \;\Longleftrightarrow\; 1 + 4\mu > 3.964
  \;\Longleftrightarrow\; \mu > 0.74 .
\end{equation}
A rigid mass would need three quarters of the plate mass. The real $\mu$ is only
$0.25$, but the plate sees $\mu_\mathrm{eff}$:
\begin{center}
\begin{tabular}{|l|l|l|}
\hline
$f$ & $\beta$ & $\mu_\mathrm{eff} = 0.25/(1-\beta)$ \\
\hline
$1031.4\,\mathrm{Hz}$ & $0.739$ & $0.96$ \\
$1046\,\mathrm{Hz}$ & $0.760$ & $1.04$ \\
$1100\,\mathrm{Hz}$ & $0.840$ & $1.57$ \\
$1190\,\mathrm{Hz}$ & $0.983$ & $15$ \\
\hline
\end{tabular}
\end{center}
All are above $0.74$. Inserting the value $1.04$ at $1046\,\mathrm{Hz}$ into
Eq.~\eqref{eq:rigid} gives
\begin{equation}
  f = \frac{2378.2}{\sqrt{1 + 4\cdot 1.04}} = 1046\,\mathrm{Hz} .
\end{equation}
The $d=4$ line of the lower-edge table \emph{is} the rigid-mass formula, with the
mass evaluated at the frequency where it acts.

\paragraph{Result.}
This is a \emph{Bragg gap with a resonance-boosted mass}:
\begin{itemize}
  \item the \textbf{upper edge} comes from the lattice (sin pattern at $X$,
        untouched by the resonator);
  \item the \textbf{lower edge} comes from the resonator's growing effective
        mass pulling band 1 at $M$ down.
\end{itemize}
It works this well because $f_r$ sits just above the $X$ fold. The width is
$164\,\mathrm{Hz}$ ($14.7\,\%$, FE), compared with $142\,\mathrm{Hz}$ ($11\,\%$)
from the one-wave estimate.

% ---------------------------------------------------------------------
\section{When is the gap above or below $f_r$?}
\label{sec:gap_regimes}

A sweep of $f_r$ with the RM plane-wave model, $\mu = 0.25$, $12$ $\mathbf{k}$-points
per path segment, gap between bands 1 and 2. Edges are good to a few Hz (this
coarser sweep gives $1032\,\mathrm{Hz}$ where the converged value is
$1031.4\,\mathrm{Hz}$).
\begin{center}
\begin{tabular}{|l|l|l|l|l|}
\hline
$f_r$ [Hz] & Gap [Hz] & Width [Hz] & Gap relative to $f_r$ & Upper edge set by \\
\hline
500 & 487--554 & 67 & contains $f_r$ & resonator band \\
800 & 748--874 & 126 & contains $f_r$ & resonator band \\
1000 & 901--1076 & 176 & contains $f_r$ & resonator band \\
1100 & 969--1170 & 201 & contains $f_r$ & resonator band \\
1150 & 1001--1194 & 193 & contains $f_r$ & sin pattern at $X$ \\
1200 & 1032--1194 & 162 & below $f_r$ & sin pattern at $X$ \\
1300 & 1090--1194 & 105 & below $f_r$ & sin pattern at $X$ \\
1400 & 1142--1194 & 52 & below $f_r$ & sin pattern at $X$ \\
1500 & 1190--1194 & 5 & below $f_r$ & sin pattern at $X$ \\
1600 & none & -- & -- & -- \\
\hline
\end{tabular}
\end{center}

\paragraph{Rule 1: this gap can never lie entirely above $f_r$.}
Band 1 always stays below $f_r$. Take $\lambda$ between $0$ and the smaller of
$\lambda_r$ and the lowest bare $\lambda_n^0$ at this $\mathbf{k}$. Then
$\mu_\mathrm{eff} > 0$ and every term of the sum in Eq.~\eqref{eq:pwe} is
positive. At $\lambda = 0$ the right side is $0$. As $\lambda$ approaches the
smaller limit, the right side goes to $+\infty$, either through
$\mu_\mathrm{eff}$ or through the lowest term. In between it rises steadily, so
it equals $1$ exactly once. That root is band 1: below $f_r$ \emph{and} below the
bare plate band, at every $\mathbf{k}$.

\paragraph{Rule 2: the upper edge is the lower of two candidates.}
\begin{itemize}
  \item the resonator band near $\Gamma$, at about
        $f_r\sqrt{1+\mu-\mu S}$, which moves with $f_r$ and lies above $f_r$ as
        long as $S < 1$ (up to $f_r \approx 1765\,\mathrm{Hz}$ here);
  \item the sin pattern at $X$, fixed at $f_X$ and independent of the resonator.
\end{itemize}
The switch happens at roughly $f_r \approx 1120\,\mathrm{Hz}$. (In the sweep the
lowest point of the resonator band lies slightly off $\Gamma$ on $\Gamma$--$X$.)

\paragraph{The three regimes.}
\begin{enumerate}
  \item \textbf{$f_r$ below $f_X$: the gap contains $f_r$.} Band 1 flattens below
        $f_r$ and band 2 starts above it. This is the classical local-resonance
        gap once $f_r$ is well below the switch.
  \item \textbf{$f_r$ above $f_X$, up to about $1510\,\mathrm{Hz}$ here: the
        whole gap is below $f_r$.} The upper edge is stuck at $f_X$, and the lower
        edge is band 1 at $M$, pulled down by the growing $\mu_\mathrm{eff}$. The
        production case $f_r = 1200\,\mathrm{Hz}$ is just past the boundary.
  \item \textbf{$f_r$ above about $1510\,\mathrm{Hz}$: no complete gap.} Band 1
        at $M$ has risen past $f_X$.
\end{enumerate}
A gap above $f_r$ could only be a higher one, between bands 2 and 3. That gap did
not open for any tuning tested ($300$--$2000\,\mathrm{Hz}$), because band 2
reaches $2378\,\mathrm{Hz}$ at $M$. The widest gap found was at
$f_r = 1100\,\mathrm{Hz}$ ($201\,\mathrm{Hz}$), close to where the upper edge
switches candidates (candidate for the stage~7 parameter study).

% ---------------------------------------------------------------------
\section{Check questions with answers}

\paragraph{Question 1.} \emph{Double the resonator mass while keeping $f_r$
($\mu = 0.5$). What happens to each gap edge, and why does one of them not move?}

\emph{Answer.} The gap widens. The larger mass drags the lower edge down
(mass loading through $\mu_\mathrm{eff}$), from $1032$ to $927\,\mathrm{Hz}$. The
upper edge stays at $1194.4\,\mathrm{Hz}$. The reason is not that it is
sinusoidal (the cos pattern is too), but that the sin pattern has a \emph{node at
the resonator}: $w_0 = 0$, so no force acts on it. Width: $162 \to 268\,\mathrm{Hz}$.

\paragraph{Question 2.} \emph{Tune the resonator to $f_r = 800\,\mathrm{Hz}$, well
below the $X$ frequency of $1194\,\mathrm{Hz}$. Which mechanisms still act, what
now determines the edges, and does the gap sit below or above $f_r$?}

\emph{Answer.} This is different from the $1200\,\mathrm{Hz}$ case. The model
gives $748$--$874\,\mathrm{Hz}$:
\begin{itemize}
  \item it is a local-resonance gap, not mostly Bragg: neither edge is a Bragg
        frequency, and the sin pattern at $X$ lies far above the gap and plays no
        role;
  \item the gap \emph{straddles} $f_r$, starting below $800\,\mathrm{Hz}$ and
        ending above it;
  \item $\mu$ still matters, and now it moves \emph{both} edges: with
        $\mu = 0.5$ the gap becomes $707$--$936\,\mathrm{Hz}$. The upper edge is
        no longer pinned.
\end{itemize}

\paragraph{Question 3.} \emph{Why can the upper edge not rise above about
$1194\,\mathrm{Hz}$, however heavy or stiff the resonator is? What would you have
to change in the cell to raise it?}

\emph{Answer.} The upper-edge mode does not see the resonator, because it has a
node there. One option is a configuration where this mode does see a resonator.
Moving the single resonator does not achieve that: the two waves at $X$ always
have one combination that vanishes at $\mathbf{x}_0$, wherever $\mathbf{x}_0$ is.
One would need a \emph{second resonator per cell} at a different position, where
that pattern is not zero. The other route is to raise the bare frequency itself:
for bending waves $f_X \propto h/a^2$, so a smaller cell or a thicker plate.
